% \changes{2.1}{2018/06/06}{Unterstütze TUM Neue Helvetica via fontspec}
% \changes{2.2}{2018/06/07}{Doku aufräumen}
% \changes{2.3}{2019/03/18}{Kleinere Verbesserungen und Aktualisierungen, Details im git}
% \changes{2.4}{2019/03/20}{Umstellung auf CC 4.0, korrigiert auch Typo im engl. Text}
% \changes{2.5}{2019/07/23}{Literaturverzeichnisse mit biblatex}
%
% \subsection{Initialisierung der Dokumentenklasse}
%
% \subsubsection{Identifizierungsabschnitt der Dokumentenklasse}
%
% \begin{macro}{\ldv@classversion}
%   Die Quelldatei kann zwei \LaTeX-Klassendateien erzeugen, eine
%   Artikel- und eine Buchklasse. Die Versionsangabe
%   ist für alle gleich.
%
%    \begin{macrocode}
\NeedsTeXFormat{LaTeX2e}
\newcommand*\ldv@classversion{2019/07/23 Dokumentenklasse des LDV -
  Version 2.5}
%<+article>\ProvidesClass{ldvarticle}[\ldv@classversion]
%<+article>\newcommand*\ldv@komaclass{scrartcl}
%<+book>\ProvidesClass{ldvbook}[\ldv@classversion]
%<+book>\newcommand*\ldv@komaclass{scrbook}
%    \end{macrocode}
% \end{macro}
%
%
%
% \subsubsection{Deklaration der Klassenoptionen}
% \label{sec:deklaration-optionen}
%
% Unsere Dokumentenklassen arbeiten mit dem Paket \emph{xkeyval}. Alle
% Optionen werden also -- anders als in den Standardklassen, aber
% ähnlich wie in vielen neuen Paketen -- durch ein Schlüssel-Wert-Paar
% vom Benutzer eingestellt.
%
% \begin{macro}{ldv@choicekeyval}
% \begin{macro}{ldv@choicekeynr}
%   xkeyval erzeugt beim Laden automatisch eine interne Kopie der
%   Optionenliste (also von \verb|\@classoptionslist|). Danach können
%   alle Optionen eingestellt werden.
%
%   Die Makros \verb|ldv@choicekeyval| und \verb|ldv@choicekeynr|
%   verwende ich später zur Auswertung aller choice
%   key-Optionen. Ersteres beinhaltet dabei den textuellen Wert und
%   letzteres den numerischen Wert, welcher gemäß der Reihenfolge in
%   der Optionen-Definition bestimmt ist.
%
%    \begin{macrocode}
\RequirePackage{xkeyval}
\newcommand*\ldv@choicekeyval\@empty
\newcommand*\ldv@choicekeynr\@empty
%    \end{macrocode}
% \end{macro}
% \end{macro}
%
%
%
% Das etoolbox-Paket stellt nützliche Hilfsfunktionen (z.\,B. if/else) zur
% Verfügung
%    \begin{macrocode}
\RequirePackage{etoolbox}
%    \end{macrocode}
%
%
% \paragraph{Im Folgenden sind die Optionen in alphabetischer Reihenfolge definiert.}
%
% \begin{option}{biblatex}
% \begin{macro}{\ldv@biblatex}
% \begin{macro}{\ifldv@useBiblatex}
%   Die Option \emph{biblatex} aktiviert die Unterstützung für moderne
%   Literaturverzeichnisse mit der Kombination biblatex und biber. Gleichzeitig
%   deaktivert sie natbib und die Neudefinitionen des \verb|\cite|-Befehls, da
%   diese mit biblatex inkompatibel und veraltet sind.
%
%   Um einen anderen Literaturverzeichnis-Stil zu wählen, kann die Klassenoption
%   \enquote{biblatex} auf einen beliebigen biblatex Stil wie
%   \enquote{authoryear-comp} gesetzt werden. Der Standardstil \enquote{ldv}
%   erweitert \emph{numeric-comp} um ein paar Eigenheiten wie Abschlussarbeiten.
%    \begin{macrocode}
\newif\ifldv@useBiblatex
\ldv@useBiblatexfalse
\DeclareOptionX<ldv>{biblatex}[ldv]{\ldv@useBiblatextrue%
  \def\ldv@biblatex{#1}}
%    \end{macrocode}
% \end{macro}
% \end{macro}
% \end{option}
%
% \begin{option}{DIV}
% \begin{macro}{\ifldv@isDivSet}
% \begin{macro}{\ldv@isDivSettrue}
% \begin{macro}{\ldv@isDivSetfalse}
%   Unterschiedliche Schriftarten benötigen unterschiedliche
%   Satzspiegel. Deshalb muss ich je nach gewählter Schriftart der
%   vorgegebenen Satzspiegel anpassen. Dies darf ich aber nur, solange
%   der Benutzer keinen eigenen DIV-Wert vorgibt. Um diesen Fall
%   erkennen zu können, muss ich die Option \enquote{DIV} überwachen. Die
%   Bedingung \verb|\ifldv@isDivSet| ist wahr, falls der Benutzer
%   einen DIV-Wert über die Klassenoptionen festgelegt hat.
%
%    \begin{macrocode}
\newif\ifldv@isDivSet
\DeclareOptionX<ldv>{DIV}{\ldv@isDivSettrue}
%    \end{macrocode}
% \end{macro}
% \end{macro}
% \end{macro}
% \end{option}
%
% \begin{option}{doctype}
% \begin{macro}{\ifldv@studthesis}
% \begin{macro}{\ldv@studthesistrue}
% \begin{macro}{\ldv@studthesisfalse}
% \begin{macro}{\ldv@doctype}
%   Die Option \verb|doctype| aktiviert die Einstellungen und Makros für
%   eine studentische Abschlussarbeit. Dazu gibt es zum einen den
%   globalen Wahrheitswert \verb|\ifldv@studthesis|. Er wird überall
%   benutzt, wo geprüft werden soll, ob gerade irgendeine studentische
%   Abschlussarbeit gesetzt wird. Zum anderen speichert das Makro
%   \verb|ldv@doctype| den konkreten Typ der Abschlussarbeit in
%   Kleinschreibung (durch das Sternchen von
%   \verb|define@choicekey*|).
%
%    \begin{macrocode}
%<*book>
\newif\ifldv@studthesis
\newif\ifldv@phdthesis
\define@choicekey*+{ldv}{doctype}[\ldv@doctype]{%
  phdthesis,diplomarbeit,mastersthesis,studienarbeit,bachelorsthesis,idp%
}{%
  \newcommand*\tempa{phdthesis}%
  \ifx\ldv@doctype\tempa \ldv@phdthesistrue\else\ldv@studthesistrue\fi%
}{%
  \ClassWarning{\@currname}{%
    The value given for the option doctype is not known.%
  }
}
%</book>
%    \end{macrocode}
% \end{macro}
% \end{macro}
% \end{macro}
% \end{macro}
% \end{option}
%
% \begin{option}{fontstyle}
% \begin{macro}{\ldv@fontstyle}
% \begin{macro}{\ldv@fontstylenr}
%   Mit der Option \emph{fontstyle} kann der Benutzer eine von drei
%   Schriftarten auswählen. Die Makros \verb|ldv@fontstyle| und
%   \verb|\ldv@fontstylenr| speichern diese, Ersteres in wörtlicher,
%   Zweiteres in numerischer Form. Die Änderung der Schriftart
%   passiert tatsächlich in Kapitel~\ref{sec:layout}; dort finden sich
%   auch die genauen Werte dieser Option. Zusätzlich hängt auch die
%   Wahl des DIV-Faktors für den Satzspiegel von dieser Option
%   ab. Diese Einstellung kann erst nach dem Verarbeiten der
%   Klassenoptionen vorgenommen werden
%   (Abschnitt~\ref{sec:impl:seitenlayout}), weil beide Optionen,
%   \enquote{fontstyle} und \enquote{DIV}, hierfür geparst sein müssen.
%
%    \begin{macrocode}
\newcommand*\ldv@fontstylenr\@empty
\define@choicekey*+{ldv}{fontstyle}[\ldv@fontstyle\ldv@fontstylenr]%
  {sans,serif}{}{%
    \ClassWarning{\@currname}{%
      The value given for the option fontstyle is not known.%
    }
  }
%    \end{macrocode}
% \end{macro}
% \end{macro}
% \end{option}
%
% \begin{option}{inputenc}
% \begin{macro}{\ldv@inputenc}
% \begin{macro}{\ldv@defaultinputenc}
%   Die Option \emph{inputenc} setzt das Makro \verb|\ldv@inputenc|
%   auf den zu verwendenden Zeichensatz für das Paket
%   \emph{inputenc}. Der Standardwert ist UTF-8
%   (\verb|\ldv@defaultinputenc|).
%
%    \begin{macrocode}
\newcommand*\ldv@defaultinputenc{utf8}
\DeclareOptionX<ldv>{inputenc}[\ldv@defaultinputenc]{\def\ldv@inputenc{#1}}
%    \end{macrocode}
% \end{macro}
% \end{macro}
% \end{option}
%
% \begin{option}{lang}
% \begin{option}{english}
% \begin{option}{ngerman}
% \begin{macro}{\ldv@lang}
%   Mit der Option \emph{lang} kann der Benutzer die Textsprache
%   einstellen. Der Vorgabewert muss als globale Klassenoption in die
%   Optionenliste eingetragen werden, damit er von anderen Paketen
%   (z.B. babel und varioref) genutzt wird. Dadurch dass der
%   Vorgabewert hier ganz am Anfang der Optionenliste steht,
%   überschreiben Sprachoptionen des Benutzers diese
%   Vorgabe. (Tatsächlich lädt Babel alle Sprachen die in der
%   Optionenliste stehen; die letzte bleibt dann aktiv.)
%
%   Diese Klassen laden deutsch und englisch. Deutsch wird im Falle
%   einer Dissertation für die Titelseite gebraucht. Englisch wird als
%   zweites geladen, so dass es die standardmäßig aktive Sprache ist.
%
%   Das Makro \verb|\ldv@lang| speichert die vom Benutzer zuletzt
%   ausgewählte Sprache, damit sie nach dem Laden des Babel-Pakets
%   aktiviert werden kann. Denn Babel aktiviert eine Sprache kein
%   zweites Mal. Wählt der Benutzer also ngerman mit
%   \verb|lang=ngerman| aus, bleibt das ohne Effekt. Durch die
%   zusätzliche Aktivierung weiter unten hat \verb|lang=ngerman| dann
%   doch den vom Benutzer gewünschten Effekt.
%
%   Die zusätzlichen Optionen \emph{ngerman} und \emph{english} sind
%   nötig, damit diese Sprachen auch mit dem neuen Mechanismus
%   behandelt werden.  Das heißt, beide Sprachen werden automatisch
%   von Babel aufgrund der hier erweiterten Liste an Klassenoptionen
%   geladen, egal welche Optionen der Benutzer eingibt.  Mit den
%   zusätzlich deklarierten Optionen ngerman und english werden
%   Benutzerangaben der Form
% \begin{verbatim}
% \documentclass[ngerman]{ldvbook}
% \documentclass[lang=ngerman,english]{ldvbook}
% \end{verbatim}
%   auch korrekt verarbeitet.  Die jeweils zuletzt angegebene Sprache
%   wird am Dokumentenanfang schließlich aktiviert.
%
%    \begin{macrocode}
\edef\@classoptionslist{ngerman,english,\@classoptionslist}
\newcommand*\ldv@lang{}
\DeclareOptionX<ldv>{lang}[english]{%
  \XKV@addtolist@n\@classoptionslist{#1}%
  \renewcommand*\ldv@lang{#1}
}
\DeclareOptionX<ldv>{english}{%
  \renewcommand*\ldv@lang{english}
}
\DeclareOptionX<ldv>{ngerman}{%
  \renewcommand*\ldv@lang{ngerman}
}
%    \end{macrocode}
% \end{macro}
% \end{option}
% \end{option}
% \end{option}
%
% \begin{option}{omitpackage}
%   Die Pakete hyperref und varioref führen manchmal zu
%   Kompatibilitäts\-problemen mit anderen Paketen. Mit der Option
%   omitpackage kann man angeben, welches Paket nicht automatisch
%   geladen werden soll. Die Option kann mehrmals erscheinen, um das
%   Laden mehrerer Pakete zu verhindern. Beispielsweise bewirkt
%\begin{verbatim}
%\documentclass[omitpackage=hyperref,omitpackage=varioref]{ldvbook}
%\end{verbatim}
%   dass die Pakete hyperref und varioref nicht automatisch geladen
%   werden.
%
%    \begin{macrocode}
\newif\ifldv@doloadhyperref
\ldv@doloadhyperreftrue
\newif\ifldv@doloadvarioref
\ldv@doloadvarioreftrue
\define@choicekey*+{ldv}{omitpackage}[\ldv@choicekeyval\ldv@choicekeynr]%
  {hyperref,varioref}{%
    \ifcase\ldv@choicekeynr\relax
      \ldv@doloadhyperreffalse
    \or
      \ldv@doloadvarioreffalse
    \fi
  }{%
    \ClassWarning{\@currname}{%
      You can only omit the packages hyperref and varioref.%
    }%
  }
%    \end{macrocode}
% \end{option}
%
% Als Letztes wird der Mechnismus gesetzt, dass alle Optionen, die von
% dieser Klassendatei nicht bearbeitet wurden, an die jeweilige
% \KOMAScript-Klasse weitergeleitet werden. Damit ist das Verhalten
% für alle Optionen konfiguriert.
%
%    \begin{macrocode}
\DeclareOptionX*{\PassOptionsToClass{\CurrentOption}{\ldv@komaclass}}
%    \end{macrocode}
%
%
% \subsubsection{Ausführung der Optionen}
%
% In diesem Abschnitt werden zuerst Standardwerte für bestimmte
% Optionen ausgeführt. Das heißt, dass die jeweils für eine bestimmte
% Option deklarierte Aktion ausgeführt wird, falls die Option nicht
% schon vom Benutzer gesetzt wurde. Schließlich verarbeitet der Befehl
% \verb|\ProcessOptionsX| die Optionenliste des Benutzers.
%
%    \begin{macrocode}
\ExecuteOptionsX<ldv>{%
  inputenc=\ldv@defaultinputenc,%
  fontstyle=sans%
}
\ProcessOptionsX<ldv>\relax
%    \end{macrocode}
%
%
%
% \subsubsection{Seitenlayout bestimmen}
% \label{sec:impl:seitenlayout}
%
% Falls die Schriftart Helvetica ausgewählt wurde und der Benutzer
% keinen DIV-Wert vorgegeben hat, stelle ich hier den Wert 9 als
% Vorgabe ein. Für die anderen Schriften behalten ich den Standardwert
% 10 der \KOMAScript-Klassen bei.
%
%    \begin{macrocode}
\ifnum\ldv@fontstylenr=0%
  \ifldv@isDivSet\else
    \PassOptionsToClass{DIV=9}{\ldv@komaclass}%
    \ClassInfo{ldvklassen}{Set DIV to 9}%
  \fi
\fi
%    \end{macrocode}
%
% Der folgende Code berechnet ähnlich wie die Option DIV=calc der
% \KOMAScript-Klassen einen Schriftart- und Blattgrößen-abhängigen
% DIV-Wert, jedoch für eine größere mittlere Zeilenbreite. Dies kann
% an dieser Stelle aber noch nicht geschehen, weil die
% Schriftgrößendateien erst von den \KOMAScript-Klassen geladen
% werden. Bisher kann ich mich nicht geeignet in die
% \KOMAScript-Klassen einklinken. Deshalb bleibt dieser Code vorerst
% ungenutzt, auch wenn er schon passende DIV-Werte berechnet.
%
% \begin{macro}{\ldv@laxLineWidth}
% \begin{macro}{\ldv@bcor}
% \begin{macro}{\ldv@div}
%   \verb|\ldv@laxLineWidth| speichert die berechnete mittlere
%   Zeilenlänge für den laxen Berechnungsalgorithmus, \verb|\ldv@bcor|
%   die vom Benutzer gewünschte Bindekorrektur und \verb|\ldv@div| den
%   berechneten oder vom Benutzer vorgegebenen DIV-Wert
%   (Teilungsfaktor für den Satzspiegel).
%
%    \begin{macrocode}
\newlength\ldv@laxLineWidth
\newlength\ldv@bcor  %% TODO
\newcount\ldv@div
%    \end{macrocode}
% \end{macro}
% \end{macro}
% \end{macro}
%
% \begin{macro}{\ldv@roundeddiv}
%   Berechnet die gerundete ganzzahlige Division zweier Dimensionen.
%
%    \begin{macrocode}
\newcommand*\ldv@roundeddiv[2]{
  \newlength\ldv@tempdima
  \setlength\ldv@tempdima{#1}
  \divide\ldv@tempdima by#2
  \ldv@div\ldv@tempdima
  \ifnum \ldv@div<4
    \ldv@div=4
  \else
    \newlength\ldv@tempdimb
    \setlength\ldv@tempdimb{#1}%
    \divide\ldv@tempdimb by\ldv@div
    \addtolength\ldv@tempdima{\ldv@tempdima}%
    \addtolength\ldv@tempdima{-\ldv@tempdimb}%
    \advance\ldv@div by\@ne
    \setlength\ldv@tempdimb{#1}%
    \divide\ldv@tempdimb by\ldv@div
    \ifdim \ldv@tempdima<\ldv@tempdimb
      \advance\ldv@div by\m@ne
    \fi
  \fi
}
%    \end{macrocode}
% \end{macro}
%
%    \begin{macrocode}
\newif\ifldv@mincl  %% TODO
\ldv@minclfalse
%    \end{macrocode}
%
% \begin{macro}{\ldv@getLaxLineWidth}
%   Der \verb|DIV=calc|-Algorithmus des typearea-Pakets berechnet ein
%   Seitenlayout für eine optimale mittlere Zeilenbreite von etwas
%   über 70 Zeichen pro Zeile (\verb|\ta@temp@goodwidth|). Das
%   entspricht zwar guter typografischer Praxis, führt aber für ein
%   DIN A4-Blatt zu sehr großen Rändern (DIV-Wert von 8 oder kleiner),
%   insbesondere bei Helvetica in 11pt.  \verb|\ldv@getLaxLineWidth|
%   berechnet dagegen eine etwas laxere mittlere Zeilenbreite von etwa
%   85 Zeichen. Bei CMR führt dies zu einem DIV-Wert von 10. Damit
%   kann die Berechnungsautomatik für den Seitenspiegel als
%   Standardeinstellung genutzt werden, weil sie kompatibel mit der
%   Standardeinstellung der \KOMAScript-Klassen ist.
%    \begin{macrocode}
\newcommand*\ldv@getLaxLineWidth{%
  \settowidth\ldv@laxLineWidth{\normalfont%
%    \end{macrocode}
% Zur Bestimmung der mittleren Zeichenlänge berücksichtige ich die
% Häufig\-keits\-verteilung der Zeichen in der deutschen und
% englischen Sprache (einschließlich Leerzeichen, beide Verteilungen
% gemäß http://en.wikipedia.org/ wiki/Letter\_frequencies);
% Buchstaben, welche hier nicht erscheinen haben, eine Häufigkeit von
% unter 1\%. Zuerst nun die Buchstabenverteilung in der deutschen
% Sprache:
%    \begin{macrocode}
eeeeeeeeeeeeeeeee nnnnnnnnnn iiiiiiii sssssss rrrrrrr aaaaaaa tttttt ddddd hhhhh uuuu lll ccc ggg mmm ooowwbbffkzv%
%    \end{macrocode}
% Dann die Verteilung in der englischen Sprache:
%    \begin{macrocode}
eeeeeeeeeeeee ttttttttt aaaaaaaa oooooooo iiiiiii nnnnnnn ssssss hhhhhh rrrrrr dddd llll ccc uuu mm ww ff gg yy pp bvk}%
%    \end{macrocode}
% Nun skaliere ich die gewonnene Breite noch so, dass ich bei der
% normalen \LaTeX-Schrift (Computer Modern Serif) einen DIV-Faktor von
% 10 erhalten (Standardeinstellung von \KOMAScript\ bei DIN A4 und
% 11pt) und bei der Helvetica den DIV-Wert 9. Damit erhalte ich in
% verschiedenen Tests 85--87 Zeichen pro Zeile.
%    \begin{macrocode}
  \setlength\ldv@laxLineWidth{.38\ldv@laxLineWidth}%
}
%    \end{macrocode}
% \end{macro}
%
% \begin{macro}{\ldv@getDIV}
%   Bestimmt einen geeigneten DIV-Faktor in Abhängigkeit von der
%   Schrift und den Papiermaßen. Das Ergebnis schreibt
%   \verb|\ldv@getDIV| dann in das Makro \verb|\ldv@div|.
%
%    \begin{macrocode}
\newcommand*\ldv@getDIV{%
%    \end{macrocode}
% Das wesentliche Problem der DIV-Wert-Berechung ist die Berechnung
% der Streifenbreite. Dazu bestimme ich am Anfang die gesamte
% Randbreite aus linkem und rechten Rand.
%    \begin{macrocode}
  \newlength\ldv@hBandWidth%
  \setlength\ldv@hBandWidth{\paperwidth}%
  \addtolength\ldv@hBandWidth{-\ldv@bcor}%
  \addtolength\ldv@hBandWidth{-\ldv@laxLineWidth}%
  \if@twocolumn%
    \addtolength\ldv@hBandWidth{-\columnsep}%
    \addtolength\ldv@hBandWidth{-\ldv@laxLineWidth}%
  \fi%
%    \end{macrocode}
% Der Rand soll aus drei Streifen bestehen.
%    \begin{macrocode}
  \divide\ldv@hBandWidth by3%
%    \end{macrocode}
% Laut dem Quellcode des typearea-Pakets schrumpft die Streifenbreite
% auf 75\%, wenn die Randnotizen in den Satzspiegel eingeschlossen
% werden sollen.
%    \begin{macrocode}
  \ifldv@mincl%
    \setlength\ldv@hBandWidth{.75\ldv@hBandWidth}%
  \fi%
%    \end{macrocode}
% Ein Streifen soll aber immer mindestens 5\,mm breit sein.
%    \begin{macrocode}
  \ifdim \ldv@hBandWidth <5mm%
    \setlength\ldv@hBandWidth{5mm}%
  \fi%
%    \end{macrocode}
% Jetzt steht die Streifenbreite abschließend fest. Daraus ergibt sich
% durch eine gerundete Ganzzahldivision der DIV-Faktor (in
% \verb|\ldv@div|).
%    \begin{macrocode}
  \ldv@roundeddiv{\paperwidth}{\ldv@hBandWidth}%
  \ClassInfo{\@currname}{Computed DIV=\the\ldv@div.}%
}
%    \end{macrocode}
% \end{macro}
%
% \begin{macro}{\ldv@autotypearea}
%   Stellt einen geeigneten Satzspiegel in Abhängigkeit der Schrift
%   und der Papiermaße ein.
%
%    \begin{macrocode}
\newcommand*\ldv@autotypearea{%
  \ldv@getLaxLineWidth%
  \ldv@getDIV%
  \PassOptionsToClass{DIV=\ldv@div}{ldv@komaclass}
}
%% Das ist hier noch nicht möglich, weil die Schriftgröße noch nicht
%% eingestellt ist.
%% \ldv@autotypearea
%    \end{macrocode}
% \end{macro}
%
%
%
% \subsubsection{Laden der Basisklasse}
%
% Nun wurden alle Optionen verarbeitet und unsere Klassendatei
% konfiguriert. Jetzt können wir die \KOMAScript-Klassendatei laden,
% auf die wir aufbauen. Dabei werden -- wie in
% Kapitel~\ref{sec:deklaration-optionen} bereits beschrieben -- die
% von uns nicht verarbeiteten Optionen an diese Klasse weitergereicht.
%
%    \begin{macrocode}
\LoadClass{\ldv@komaclass}
%    \end{macrocode}
%
%
% \subsection{Kodierung und Sprache}
%
% In dieser Klassendatei laden wir viele Pakete, um dem Benutzer das
% Leben einfacher zu machen. Er soll eine komfortable Umgebung einfach
% nur mit der Zeile
%
% \verb|\documentclass{ldvarticle}|
%
% \noindent erhalten. Deshalb werden in den folgenden Kapiteln alle
% Pakete geladen, die wir als \enquote{Standard} ansehen.
%
% Dazu gehören alle \enquote{Sprachpakete}, also Zeichenkodierung,
% Schriftenkodierung und \enquote{Kodierung der menschlichen Sprache}. Der
% Vorgabewert für die Zeichenkodierung ist UTF-8 (siehe
% Kapitel~\ref{sec:deklaration-optionen}), für die Sprache
% Englisch. Über die Klassenoptionen kann der Autor andere Werte
% einstellen. Die Übergabe der Sprache an das babel-Paket erfolgt über
% die globalen Klassenoptionen. So wirkt sie sich auch auf das
% varioref-Paket aus (siehe Kapitel~\ref{sec:verweise}).
%
% Der Aufruf von \verb|\main@language| nach dem Laden des Babel-Pakets
% stellt die Standardsprache ein, welche bei \verb|\begin{document}|
% geladen wird. Dies ist der in der Babel-Dokumentation empfohlene
% Weg für Autoren von Sprachpaketen. \verb|\ldv@lang| enthält die
% Sprache, welche vom Benutzer zuletzt angegeben wurde.
%
%    \begin{macrocode}
\RequirePackage{silence}
\WarningFilter{inputenc}{inputenc package ignored with utf8 based engines}
\RequirePackage[\ldv@inputenc]{inputenc}
\RequirePackage[T1]{fontenc}
\ifx\ldv@lang\@empty
  \RequirePackage{babel}
\else
  \RequirePackage[main=\ldv@lang]{babel}
  \expandafter\main@language\expandafter{\ldv@lang}
\fi
\RequirePackage{csquotes}
%    \end{macrocode}
%
%
%
%
% \subsection{Strukturbefehle für Fließtext}
%
% Bei der Textauszeichnung benötigen Autoren häufig verbesserte
% Funktionen und mehr Flexibilität im Bereich der Tabellen und der
% Darstellung von vorformatierten Text (z.B. Quellcode). Dazu laden
% wir die Standardpakete \emph{array} und \emph{fancyvrb}.
%
%    \begin{macrocode}
\RequirePackage{array}
\RequirePackage{fancyvrb}
\fvset{xleftmargin=2em}
%    \end{macrocode}
%
%
% \begin{macro}{\emphemph}
%   \LaTeX\ besitzt ein Makro, um Text einfach hervorzuheben. Häufig
%   ist aber ein zweistufige Betonung erwünscht -- ähnlich dem
%   Textmodell von HTML. Beispielsweise könnte \verb|\emph| den Text
%   kursiv darstellen und \verb|\emphemph| -- die zweite Stufe -- in
%   blauer Farbe. Für den Schwarz-Weiß-Druck ist die Vorgabe hier
%   Fettdruck.
%
%    \begin{macrocode}
\newcommand\emphemph[1]{{\bfseries#1}}
%    \end{macrocode}
% \end{macro}
%
%
% \begin{macro}{\simpleverb}
%   Der Befehl zum Darstellen von direkt formatierten Text,
%   \verb|\verb|, kann nicht in Überschrift, Bildunterschriften und
%   ähnlichem verwendet werden. Deshalb braucht man einen
%   Ersatzbefehl, der zwar nicht das direkte Verhalten nachbilden
%   kann, aber zumindest das Ausssehen des Textes.
%
%    \begin{macrocode}
\newcommand\simpleverb[1]{{\ttfamily#1}}
%    \end{macrocode}
% \end{macro}
%
%
% \begin{environment}{note}
% \begin{macro}{\notename}
%   Die \verb|note|-Umgebung kann dazu verwendet werden, eine
%   abgesetzte Passage/Bemerkung beispielsweise am Ende eines
%   Abschnitts abzudrucken. Der Name dieses Abschnitts kann durch den
%   optionalen Parameter bestimmt werden. Der Vorgabewert \enquote{Note} ist
%   sprachabhängig durch Babel realisiert.
%
%    \begin{macrocode}
\newcommand*\notename{Note}
\addto\captionsngerman{%
  \renewcommand*\notename{Bemerkung}%
}
\addto\captionsgerman{%
  \renewcommand*\notename{Bemerkung}%
}
\addto\captionsenglish{%
  \renewcommand*\notename{Note}%
}
\newenvironment{note}[1][\notename]{%
  \vspace{\baselineskip}%
  \noindent\small{\usekomafont{sectioning}#1:}\hspace{.5em}%
}{}
%    \end{macrocode}
% \end{macro}
% \end{environment}
%
%
%
%
% \subsection{Grafik}
%
% Zuerst werden die Standardpakete für den Umgang mit Farben und
% Bildern geladen. Sie werden von vielen Autoren genutzt und stellen
% heutzutage eine Basisfunktionalität in der Textverarbeitung dar.
%
%    \begin{macrocode}
\RequirePackage{xcolor}
\RequirePackage{graphicx}
\RequirePackage{tikz}
%    \end{macrocode}
%
% \begin{macro}{\graphicswidth}
% \begin{macro}{\graphicswidthtwo}
%   Alle fließenden Abbildungen eines Typs sollten in einem Dokument
%   immer die gleiche Breite haben. Dazu stellen diese
%   Dokumentenklassen für die beiden üblichen Konstellationen
%   Längen-Makros zur Verfügung: \verb|\graphicswidth| ist die
%   passende Breite, wenn eine Abbildung in der Fließumgebung ist
%   (also über die ganze Breite abzüglich eines kleinen Rands von
%   2em). Beinhaltet die Fließumgebung zwei Grafiken nebeneinander, so
%   sollen sie die Breite \verb|\graphicswidthtwo| erhalten und einen
%   Zwischenabstand von 2em.
%
%    \begin{macrocode}
\newlength\graphicswidth
\setlength\graphicswidth{\columnwidth}
\addtolength\graphicswidth{-4\parindent}
\newlength\graphicswidthtwo
\setlength\graphicswidthtwo{.5\columnwidth}
\addtolength\graphicswidthtwo{-3\parindent}
%    \end{macrocode}
% \end{macro}
% \end{macro}
%
% Bildunterschriften sollen gleich breit wie Abbildungen
% sein. Zusätzlich formatiere ich sie noch etwas um: Hervorgehobene
% Bildbezeichnungen, aber ohne Einzug.
%
%    \begin{macrocode}
%\setcapwidth[c]{\graphicswidth}
\setcapindent{0pt}
\setkomafont{caption}{\small}
\setkomafont{captionlabel}{\usekomafont{sectioning}}
%    \end{macrocode}
%
% \begin{macro}{TUMBlau}
% \begin{macro}{TUMBlau1}
% \begin{macro}{TUMBlau2}
% \begin{macro}{TUMBlau3}
% \begin{macro}{TUMBlau4}
% \begin{macro}{TUMBlau5}
% \begin{macro}{TUMDunkelgrau}
% \begin{macro}{TUMMittelgrau}
% \begin{macro}{TUMHellgrau}
% \begin{macro}{TUMGruen}
% \begin{macro}{TUMOrange}
% \begin{macro}{TUMElfenbein}
%   Die Farben der TUM Corporate Identity können als Hintergrund von
%   Boxen oder als Textumrandung verwendet werden. Sie sind hier
%   vordefiniert mit den Namen, die der Corporate Identity Reference
%   Guide nennt.
%
%    \begin{macrocode}
\definecolor{TUMBlau}       {cmyk}{1.00,0.43,0.00,0.00}
\definecolor{TUMBlau1}      {cmyk}{1.00,0.57,0.12,0.70}
\definecolor{TUMBlau2}      {cmyk}{1.00,0.54,0.04,0.19}
\definecolor{TUMBlau3}      {cmyk}{0.90,0.48,0.00,0.00}
\definecolor{TUMBlau4}      {cmyk}{0.65,0.19,0.01,0.04}
\definecolor{TUMBlau5}      {cmyk}{0.42,0.09,0.00,0.00}
\definecolor{TUMDunkelgrau} {cmyk}{0.00,0.00,0.00,0.80}
\definecolor{TUMMittelgrau} {cmyk}{0.00,0.00,0.00,0.50}
\definecolor{TUMHellgrau}   {cmyk}{0.00,0.00,0.00,0.20}
\definecolor{TUMGruen}      {cmyk}{0.35,0.00,1.00,0.20}
\definecolor{TUMOrange}     {cmyk}{0.00,0.65,0.95,0.00}
\definecolor{TUMElfenbein}  {cmyk}{0.03,0.04,0.14,0.08}
%    \end{macrocode}
% \end{macro}
% \end{macro}
% \end{macro}
% \end{macro}
% \end{macro}
% \end{macro}
% \end{macro}
% \end{macro}
% \end{macro}
% \end{macro}
% \end{macro}
% \end{macro}
%
%
%
%
%
% \subsection{Layout}
% \label{sec:layout}
%
% \begin{macro}{\ldv@setfontstyle}
%   Falls \LuaTeX{} oder \XeTeX{} genutzt werden, werden die
%   System-Schriftarten geladen, das ermöglicht den Einsatz der TUM
%   Neue Helvetica. Als Monospace-Schrift wird in diesem Fall bei
%   Verfügbarkeit die DejaVu Sans Mono genutzt, sie fügt sich sehr
%   gut ins Schriftbild ein.
%   Mit dem Makro \verb|\ldv@setfontstyle| kann der Benutzer eine von
%   zwei Schriftkombinationen auswählen:
%   \begin{description}
%   \item[sans] (Vorgabe) Serifenlose Schrift (Helvetica) für das gesamte Dokument.
%   \item[serif] Als Hauptschriftart wird eine Times-Variante genutzt.
%   \end{description}
%
%
%    \begin{macrocode}
\RequirePackage{expl3}
\ExplSyntaxOn
\sys_if_engine_luatex:TF
{
  \RequirePackage{fontspec}
  \defaultfontfeatures{Scale=MatchLowercase}
  \defaultfontfeatures[TUM Neue Helvetica]{
    UprightFont	= TUM Neue Helvetica 55 Regular,
    BoldFont	= TUM Neue Helvetica 75 Bold,
    ItalicFont	= TUM Neue Helvetica 56 Italic,
    BoldItalicFont	= TUM Neue Helvetica 76 Bold Italic,
    Ligatures	= TeX,
    Scale = 0.92,
  }
  \IfFontExistsTF{Times}{\setmainfont{Times}}{}
  \IfFontExistsTF{Times New Roman}{\setmainfont{Times New Roman}}{}
  \IfFontExistsTF{Arial}{\setsansfont{Arial}}{}
  \IfFontExistsTF{Helvetica}{\setsansfont{Helvetica}}{}
  \IfFontExistsTF{TUM Neue Helvetica}{\setsansfont{TUM Neue Helvetica}}{}
  \IfFontExistsTF{DejaVu Sans Mono}{\setmonofont{DejaVu Sans Mono}}{}
}
{
    \sys_if_engine_xetex:TF
    {
      \RequirePackage{fontspec}
      \defaultfontfeatures{Scale=MatchLowercase}
      \defaultfontfeatures[TUM Neue Helvetica]{
        UprightFont	= TUM Neue Helvetica 55 Regular,
        BoldFont	= TUM Neue Helvetica 75 Bold,
        ItalicFont	= TUM Neue Helvetica 56 Italic,
        BoldItalicFont	= TUM Neue Helvetica 76 Bold Italic,
        Ligatures	= TeX,
        Scale = 0.92,
      }
      \IfFontExistsTF{Times}{\setmainfont{Times}}{}
      \IfFontExistsTF{Times New Roman}{\setmainfont{Times New Roman}}{}
      \IfFontExistsTF{Arial}{\setsansfont{Arial}}{}
      \IfFontExistsTF{Helvetica}{\setsansfont{Helvetica}}{}
      \IfFontExistsTF{TUM Neue Helvetica}{\setsansfont{TUM Neue Helvetica}}{}
      \IfFontExistsTF{DejaVu Sans Mono}{\setmonofont{DejaVu Sans Mono}}{}
    }
    {
        \RequirePackage{courier}
        \renewcommand{\rmdefault}{ptm}
        \RequirePackage[scaled=0.92]{helvet}
    }
}
\ExplSyntaxOff
\newcommand*\ldv@setfontstyle{%
  \ifcase\ldv@fontstylenr\relax%
    \renewcommand\familydefault{\sfdefault}%
  \or%
    \renewcommand\familydefault{\rmdefault}%
  \fi%
  % The following font definitions are taken from the KOMA-Script
  % classes. Only the '\sffamily' command is left out.
  \setkomafont{disposition}{\normalcolor\bfseries}%
  \setkomafont{descriptionlabel}{\normalcolor\itshape}%
  \setkomafont{dictum}{\normalfont\normalcolor\small}%
}
%    \end{macrocode}
% \end{macro}
%
% Mit obigem Befehl kann die Schrift gemäß den Klassenoptionen
% eingestellt werden.
%
%    \begin{macrocode}
\ldv@setfontstyle
%    \end{macrocode}
%
% Erfahrungsgemäß muss in vielen Dokumenten der \verb|\tolerance|-Wert
% angepasst werden, um übervolle horizontale Boxen zu vermeiden. Der
% Standardwert von \LaTeX\ ist etwas zu streng eingestellt. Im
% Folgenden benutze ich die Einstellungen von Axel Reichert
% (http://groups.google.com/group/de.comp.text.tex/msg/c375ef11e78e7bfa).
%
%    \begin{macrocode}
\tolerance=1414
\hbadness=1414
\emergencystretch=1.5em
\hfuzz=0.3pt
\widowpenalty=10000
\vfuzz\hfuzz
\raggedbottom
%    \end{macrocode}
%
%
%
%
% \subsection{Mathematik}
%
% Die AMS-Pakete für den Mathematik-Satz sind mittlerweile der
% absolute Standard unter \LaTeX. Das Basispaket \emph{amsmath} lädt
% einige weitere spezielle AMS-Pakete. Lediglich \emph{amssymb}
% benötigen wir noch für eine größere Vielfalt an Symbolen.
%
%    \begin{macrocode}
\RequirePackage{amsmath}
\RequirePackage{amssymb}
%    \end{macrocode}
%
%
% \begin{environment}{definition}
% \begin{environment}{theorem}
% \begin{environment}{proof}
% \begin{macro}{\definitionname}
% \begin{macro}{\theoremname}
%   Standard-\LaTeX\ bietet zwar die \verb|\theorem|-Umgebung, aber
%   die genügt häufig nicht und ist etwas unflexibel. Mit den Mitteln
%   des AMS-Pakets können wir leicht jeweils eine Umgebung für die
%   mathematischen Standardsätze anbieten: Der Definition, dem Satz
%   und dem zugehörigen Beweis. Alle drei sind für Englisch und
%   Deutsch lokalisiert.
%
%    \begin{macrocode}
\newcommand*\definitionname{Definition}
\newcommand*\theoremname{Theorem}
\addto\captionsngerman{%
  \renewcommand*\definitionname{Definition}%
  \renewcommand*\theoremname{Satz}%
}
\addto\captionsgerman{%
  \renewcommand*\definitionname{Definition}%
  \renewcommand*\theoremname{Satz}%
}
\addto\captionsenglish{%
  \renewcommand*\definitionname{Definition}%
  \renewcommand*\theoremname{Theorem}%
}
\newtheorem{definition}{\definitionname}
\newtheorem{theorem}{\theoremname}
\newtheorem{proof}{\proofname}
%    \end{macrocode}
% \end{macro}
% \end{macro}
% \end{environment}
% \end{environment}
% \end{environment}
%
%
%
%
% \subsection{Verweise}
% \label{sec:verweise}
%
% Querverbindungen aufzuzeigen ist in der wissenschaftlichen Arbeit
% sehr wichtig. Deshalb sollen diese Dokumentenklassen die
% verschiedenen Formen von Online- und Offline-Verweisen gut
% unterstützen.
%
% Moderne Literaturverweise ermöglicht das Paket \emph{natbib}
% (Natural Sciences Citations und References). Es ist weitreichend
% konfigurierbar und bietet neben den üblichen numerischen
% Verweisformen auch die Autor-Jahr-Form. Natbib muss vor dem
% hyperref-Paket geladen werden.
%
% Als \BibTeX-Stil verwenden diese Dokumentenklassen den eigenen Stil
% \emph{ldv} als Voreinstellung. Er unterstützt neben den üblichen
% Feldern auch ISBN-, ISSN-, DOI- und URL-Angaben. Mit dem
% \verb|language|-Feld kann man die Silbentrennung einer anderen
% Sprache für den Titel wählen. Die Kombination aus dem
% ldv-\BibTeX-Stil und den LDV-Dokumentenklassen setzt darüber hinaus
% das Literaturverzeichnis automatisch in der aktiven
% Dokumentensprache, indem sie sich in den Babel-Mechanismus
% integrieren.
%
% Die folgenden Codeblöcke integrieren natbib in diese
% Dokumentenklassen, passen das Paket an die eigenen Wünsche an,
% setzen automatisch den \BibTeX-Stil und integrieren die Sprache des
% Literaturverzeichnisses in den babel-Mechanismus.
%
%    \begin{macrocode}
\ifldv@useBiblatex
  \RequirePackage[style=\ldv@biblatex]{biblatex}
\else
  \RequirePackage{natbib}
\fi
%    \end{macrocode}
%
% \begin{macro}{\cite}
%   Weil die Klammerung der Verweise beim natbib-Paket anders ist als
%   in Standard-\LaTeX, definiere ich sie in eine kompatible Form um:
%
%    \begin{macrocode}
\ifldv@useBiblatex
\else
  \renewcommand*\cite{\citep}
\fi
%    \end{macrocode}
% \end{macro}
%
% \begin{macro}{\ldv@latex@bibliographystyle}
% \begin{macro}{\ifldv@bibstyleset}
% \begin{macro}{\ldv@bibstylesetfalse}
% \begin{macro}{\ldv@bibstylesettrue}
% \begin{macro}{\ldv@latex@bibliography}
%   Diese Klassendateien sollen ein sinnvolles Vorgabedesign
%   anbieten. Dazu gehört auch der ldv-Bibliographiestil. Dieser Stil
%   soll aber jederzeit vom Benutzer überschrieben werden können. Weil
%   der Befehl \verb|\bibliographystyle| nur einmal aufgerufen werden
%   darf, kann der Benutzer einen vorherigen Aufruf in den
%   Dokumentenklassen nicht einfach überschreiben. Deshalb prüft der
%   erneuerte Befehl \verb|\bibliography| mittels
%   \verb|\ifldv@bibstyleset|, ob der Benutzer bereits einen
%   Bibliographiestil gesetzt hat. Falls nicht, stellt er ldv.bst als
%   \BibTeX-Stil ein. Der veränderte Befehl \verb|\bibliographystyle|
%   setzt beim gleich beim ersten Aufruf \verb|\ifldv@bibstyleset| auf
%   wahr.
%
%    \begin{macrocode}
% This if has to be defined in the outer scope, as TeX token skipping
% is unintuitivly compared to usual programming languages
\newif\ifldv@bibstyleset
\ifldv@useBiblatex
\else{%
  \let\ldv@latex@bibliographystyle=\bibliographystyle
  \renewcommand*\bibliographystyle[1]{%
    \ldv@latex@bibliographystyle{#1}\ldv@bibstylesettrue%
  }
  \let\ldv@latex@bibliography=\bibliography
  \renewcommand*\bibliography[1]{%
    \ifldv@bibstyleset\else\bibliographystyle{ldv}\fi%
    \ldv@latex@bibliography{#1}%
  }
}
\fi
%    \end{macrocode}
% \end{macro}
% \end{macro}
% \end{macro}
% \end{macro}
% \end{macro}
%
% \begin{macro}{bibstyle@ldvplain}
%   In natbib kann man zu jedem Bibliographiestil auch einen
%   Zitierstil angeben. Für unseren Bibliographiestil soll es zwei
%   Zitierstile geben. Standardmäßig soll \LaTeX das Autor-Jahr-Schema
%   verwenden (Zitierstil \emph{ldv}). Wahlweise kann der Benutzer
%   aber auch das numerische Schema \emph{ldvplain} einstellen. natbib
%   bietet dafür den Befehl \verb|citestyle}| an.
%
%    \begin{macrocode}
\newcommand\bibstyle@ldv{%
  \setcitestyle{authoryear,round,comma,aysep={,}, yysep={,},notesep={, }}%
}
\newcommand\bibstyle@ldvplain{%
  \setcitestyle{numbers,square,comma,aysep={,}, yysep={,},notesep={, }}%
  \renewcommand*\bibnumfmt[1]{##1.}%
}
%    \end{macrocode}
% \end{macro}
%
% \begin{macro}{\bbland}
% \begin{macro}{\bblandsep}
% \begin{macro}{\bblandsepauthor}
% \begin{macro}{\bbletal}
% \begin{macro}{\bbleditors}
% \begin{macro}{\bbleds}
% \begin{macro}{\bbleditor}
% \begin{macro}{\bbled}
% \begin{macro}{\bbledby}
% \begin{macro}{\bbledition}
% \begin{macro}{\bbledn}
% \begin{macro}{\bblvolume}
% \begin{macro}{\bblvol}
% \begin{macro}{\bblof}
% \begin{macro}{\bblnumber}
% \begin{macro}{\bblno}
% \begin{macro}{\bblin}
% \begin{macro}{\bblpages}
% \begin{macro}{\bblpp}
% \begin{macro}{\bblpage}
% \begin{macro}{\bblp}
% \begin{macro}{\bbleidpp}
% \begin{macro}{\bblchapter}
% \begin{macro}{\bblchap}
% \begin{macro}{\bbltechreport}
% \begin{macro}{\bbltechrep}
% \begin{macro}{\bblmthesis}
% \begin{macro}{\bblphdthesis}
% \begin{macro}{\bblfirst}
% \begin{macro}{\bblfirsto}
% \begin{macro}{\bblsecond}
% \begin{macro}{\bblsecondo}
% \begin{macro}{\bblthird}
% \begin{macro}{\bblthirdo}
% \begin{macro}{\bblfourth}
% \begin{macro}{\bblfourtho}
% \begin{macro}{\bblfifth}
% \begin{macro}{\bblfiftho}
% \begin{macro}{\bblst}
% \begin{macro}{\bblnd}
% \begin{macro}{\bblrd}
% \begin{macro}{\bblth}
% \begin{macro}{\bbljan}
% \begin{macro}{\bblfeb}
% \begin{macro}{\bblmar}
% \begin{macro}{\bblapr}
% \begin{macro}{\bblmay}
% \begin{macro}{\bbljun}
% \begin{macro}{\bbljul}
% \begin{macro}{\bblaug}
% \begin{macro}{\bblsep}
% \begin{macro}{\bbloct}
% \begin{macro}{\bblnov}
% \begin{macro}{\bbldec}
%   Das Literaturverzeichnis ist mit den folgenden Makros gesetzt. Sie
%   müssen mittels des babel-Mechnismus auf Begriffe in der aktuellen
%   Sprache gesetzt werden. Dies ermöglichen die folgenden Zeilen.
%
%    \begin{macrocode}
\newcommand*\bbland{and}
\newcommand*\bblandsep{,}
\newcommand*\bblandsepauthor{,}
\newcommand*\bbletal{et~al.}
\newcommand*\bbleditors{editors}
\newcommand*\bbleds{eds.}
\newcommand*\bbleditor{editor}
\newcommand*\bbled{ed.}
\newcommand*\bbledby{edited by}
\newcommand*\bbledition{edition}
\newcommand*\bbledn{edn.}
\newcommand*\bblvolume{volume}
\newcommand*\bblvol{vol.}
\newcommand*\bblof{of}
\newcommand*\bblnumber{number}
\newcommand*\bblno{no.}
\newcommand*\bblin{in}
\newcommand*\bblpages{pages}
\newcommand*\bblpp{pp.}
\newcommand*\bblpage{page}
\newcommand*\bblp{p.}
\newcommand*\bbleidpp{pages}
\newcommand*\bblchapter{chapter}
\newcommand*\bblchap{chap.}
\newcommand*\bbltechreport{Technical Report}
\newcommand*\bbltechrep{Tech. Rep.}
\newcommand*\bblmthesis{Master's thesis}
\newcommand*\bblphdthesis{Ph.D. thesis}
\newcommand*\bblfirst{First}
\newcommand*\bblfirsto{1st}
\newcommand*\bblsecond{Second}
\newcommand*\bblsecondo{2nd}
\newcommand*\bblthird{Third}
\newcommand*\bblthirdo{3rd}
\newcommand*\bblfourth{Fourth}
\newcommand*\bblfourtho{4th}
\newcommand*\bblfifth{Fifth}
\newcommand*\bblfiftho{5th}
\newcommand*\bblst{st}
\newcommand*\bblnd{nd}
\newcommand*\bblrd{rd}
\newcommand*\bblth{th}
\newcommand*\bbljan{January}
\newcommand*\bblfeb{February}
\newcommand*\bblmar{March}
\newcommand*\bblapr{April}
\newcommand*\bblmay{May}
\newcommand*\bbljun{June}
\newcommand*\bbljul{July}
\newcommand*\bblaug{August}
\newcommand*\bblsep{September}
\newcommand*\bbloct{October}
\newcommand*\bblnov{November}
\newcommand*\bbldec{December}
\addto\captionsngerman{%
  \renewcommand*\bbland{und}
  \renewcommand*\bblandsep{}
  \renewcommand*\bblandsepauthor{}
  \renewcommand*\bbletal{et~al.}
  \renewcommand*\bbleditors{Herausgeber}
  \renewcommand*\bbleds{Hrsg.}
  \renewcommand*\bbleditor{Herausgeber}
  \renewcommand*\bbled{Hrsg.}
  \renewcommand*\bbledby{herausgegeben von}
  \renewcommand*\bbledition{Auf\/lage}
  \renewcommand*\bbledn{Aufl.}
  \renewcommand*\bblvolume{Band}
  \renewcommand*\bblvol{Bd.}
  \renewcommand*\bblof{von}
  \renewcommand*\bblnumber{Nummer}
  \renewcommand*\bblno{Nr.}
  \renewcommand*\bblin{in}
  \renewcommand*\bblpages{Seiten}
  \renewcommand*\bblpp{S.}
  \renewcommand*\bblpage{Seite}
  \renewcommand*\bblp{S.}
  \renewcommand*\bbleidpp{Seiten}
  \renewcommand*\bblchapter{Kapitel}
  \renewcommand*\bblchap{Kap.}
  \renewcommand*\bbltechreport{Technischer Bericht}
  \renewcommand*\bbltechrep{Techn. Ber.}
  \renewcommand*\bblmthesis{Masterarbeit}
  \renewcommand*\bblphdthesis{Dissertation}
  \renewcommand*\bblfirst{Erste}
  \renewcommand*\bblfirsto{1.}
  \renewcommand*\bblsecond{Zweite}
  \renewcommand*\bblsecondo{2.}
  \renewcommand*\bblthird{Dritte}
  \renewcommand*\bblthirdo{3.}
  \renewcommand*\bblfourth{Vierte}
  \renewcommand*\bblfourtho{4.}
  \renewcommand*\bblfifth{F\^^b{u}nfte}
  \renewcommand*\bblfiftho{5.}
  \renewcommand*\bblst{.}
  \renewcommand*\bblnd{.}
  \renewcommand*\bblrd{.}
  \renewcommand*\bblth{.}
  \renewcommand*\bbljan{Januar}
  \renewcommand*\bblfeb{Februar}
  \renewcommand*\bblmar{M\^^b{a}rz}
  \renewcommand*\bblapr{April}
  \renewcommand*\bblmay{Mai}
  \renewcommand*\bbljun{Juni}
  \renewcommand*\bbljul{Juli}
  \renewcommand*\bblaug{August}
  \renewcommand*\bblsep{September}
  \renewcommand*\bbloct{Oktober}
  \renewcommand*\bblnov{November}
  \renewcommand*\bbldec{Dezember}
}
\addto\captionsgerman{%
  \renewcommand*\bbland{und}
  \renewcommand*\bblandsep{}
  \renewcommand*\bblandsepauthor{}
  \renewcommand*\bbletal{et~al.}
  \renewcommand*\bbleditors{Herausgeber}
  \renewcommand*\bbleds{Hrsg.}
  \renewcommand*\bbleditor{Herausgeber}
  \renewcommand*\bbled{Hrsg.}
  \renewcommand*\bbledby{herausgegeben von}
  \renewcommand*\bbledition{Auf\/lage}
  \renewcommand*\bbledn{Aufl.}
  \renewcommand*\bblvolume{Band}
  \renewcommand*\bblvol{Bd.}
  \renewcommand*\bblof{von}
  \renewcommand*\bblnumber{Nummer}
  \renewcommand*\bblno{Nr.}
  \renewcommand*\bblin{in}
  \renewcommand*\bblpages{Seiten}
  \renewcommand*\bblpp{S.}
  \renewcommand*\bblpage{Seite}
  \renewcommand*\bblp{S.}
  \renewcommand*\bbleidpp{Seiten}
  \renewcommand*\bblchapter{Kapitel}
  \renewcommand*\bblchap{Kap.}
  \renewcommand*\bbltechreport{Technischer Bericht}
  \renewcommand*\bbltechrep{Techn. Ber.}
  \renewcommand*\bblmthesis{Masterarbeit}
  \renewcommand*\bblphdthesis{Dissertation}
  \renewcommand*\bblfirst{Erste}
  \renewcommand*\bblfirsto{1.}
  \renewcommand*\bblsecond{Zweite}
  \renewcommand*\bblsecondo{2.}
  \renewcommand*\bblthird{Dritte}
  \renewcommand*\bblthirdo{3.}
  \renewcommand*\bblfourth{Vierte}
  \renewcommand*\bblfourtho{4.}
  \renewcommand*\bblfifth{F\^^b{u}nfte}
  \renewcommand*\bblfiftho{5.}
  \renewcommand*\bblst{.}
  \renewcommand*\bblnd{.}
  \renewcommand*\bblrd{.}
  \renewcommand*\bblth{.}
  \renewcommand*\bbljan{Januar}
  \renewcommand*\bblfeb{Februar}
  \renewcommand*\bblmar{M\^^b{a}rz}
  \renewcommand*\bblapr{April}
  \renewcommand*\bblmay{Mai}
  \renewcommand*\bbljun{Juni}
  \renewcommand*\bbljul{Juli}
  \renewcommand*\bblaug{August}
  \renewcommand*\bblsep{September}
  \renewcommand*\bbloct{Oktober}
  \renewcommand*\bblnov{November}
  \renewcommand*\bbldec{Dezember}
}
\addto\captionsenglish{%
  \renewcommand*\bbland{and}
  \renewcommand*\bblandsep{,}
  \renewcommand*\bblandsepauthor{,}
  \renewcommand*\bbletal{et~al.}
  \renewcommand*\bbleditors{editors}
  \renewcommand*\bbleds{eds.}
  \renewcommand*\bbleditor{editor}
  \renewcommand*\bbled{ed.}
  \renewcommand*\bbledby{edited by}
  \renewcommand*\bbledition{edition}
  \renewcommand*\bbledn{edn.}
  \renewcommand*\bblvolume{volume}
  \renewcommand*\bblvol{vol.}
  \renewcommand*\bblof{of}
  \renewcommand*\bblnumber{number}
  \renewcommand*\bblno{no.}
  \renewcommand*\bblin{in}
  \renewcommand*\bblpages{pages}
  \renewcommand*\bblpp{pp.}
  \renewcommand*\bblpage{page}
  \renewcommand*\bblp{p.}
  \renewcommand*\bbleidpp{pages}
  \renewcommand*\bblchapter{chapter}
  \renewcommand*\bblchap{chap.}
  \renewcommand*\bbltechreport{Technical Report}
  \renewcommand*\bbltechrep{Tech. Rep.}
  \renewcommand*\bblmthesis{Master's thesis}
  \renewcommand*\bblphdthesis{Ph.D. thesis}
  \renewcommand*\bblfirst{First}
  \renewcommand*\bblfirsto{1st}
  \renewcommand*\bblsecond{Second}
  \renewcommand*\bblsecondo{2nd}
  \renewcommand*\bblthird{Third}
  \renewcommand*\bblthirdo{3rd}
  \renewcommand*\bblfourth{Fourth}
  \renewcommand*\bblfourtho{4th}
  \renewcommand*\bblfifth{Fifth}
  \renewcommand*\bblfiftho{5th}
  \renewcommand*\bblst{st}
  \renewcommand*\bblnd{nd}
  \renewcommand*\bblrd{rd}
  \renewcommand*\bblth{th}
  \renewcommand*\bbljan{January}
  \renewcommand*\bblfeb{February}
  \renewcommand*\bblmar{March}
  \renewcommand*\bblapr{April}
  \renewcommand*\bblmay{May}
  \renewcommand*\bbljun{June}
  \renewcommand*\bbljul{July}
  \renewcommand*\bblaug{August}
  \renewcommand*\bblsep{September}
  \renewcommand*\bbloct{October}
  \renewcommand*\bblnov{November}
  \renewcommand*\bbldec{December}
}
%    \end{macrocode}
% \end{macro}
% \end{macro}
% \end{macro}
% \end{macro}
% \end{macro}
% \end{macro}
% \end{macro}
% \end{macro}
% \end{macro}
% \end{macro}
% \end{macro}
% \end{macro}
% \end{macro}
% \end{macro}
% \end{macro}
% \end{macro}
% \end{macro}
% \end{macro}
% \end{macro}
% \end{macro}
% \end{macro}
% \end{macro}
% \end{macro}
% \end{macro}
% \end{macro}
% \end{macro}
% \end{macro}
% \end{macro}
% \end{macro}
% \end{macro}
% \end{macro}
% \end{macro}
% \end{macro}
% \end{macro}
% \end{macro}
% \end{macro}
% \end{macro}
% \end{macro}
% \end{macro}
% \end{macro}
% \end{macro}
% \end{macro}
% \end{macro}
% \end{macro}
% \end{macro}
% \end{macro}
% \end{macro}
% \end{macro}
% \end{macro}
% \end{macro}
% \end{macro}
% \end{macro}
% \end{macro}
% \end{macro}
%
% Das eher unbekannte Paket \emph{varioref} bietet vor allem den
% Befehl \verb|\vref|, der wie \verb|\ref| benutzt wird und einen
% intelligenten Verweis erzeugt -- zum Beispiel \enquote{Abbildung 4.1 auf
% der vorherigen Seite}.
%
% \emph{hyperref} und \emph{varioref} vertragen sich jedoch nicht
% vollständig. Das beste Vorgehen soll laut der Dokumentation zu
% hyperref sein, dass man \emph{nameref}, welches zu hyperref gehört,
% bereits vor varioref lädt. Aber selbst dann sollen noch vereinzelte
% Ungereimtheiten verbleiben. Die Optionen zum hyperref-Paket unten
% unterdrücken die Umrandung um Link-Text.
%
%    \begin{macrocode}
\ifldv@doloadhyperref%
\RequirePackage{nameref}
\fi
\ifldv@doloadvarioref%
\RequirePackage{varioref}
\fi
\ifldv@doloadhyperref%
\RequirePackage[unicode,pdfborder={0 0 0}]{hyperref}%
\expandafter\def\expandafter\UrlBreaks\expandafter{\UrlBreaks%  save the current one
  \do\a\do\b\do\c\do\d\do\e\do\f\do\g\do\h\do\i\do\j%
  \do\k\do\l\do\m\do\n\do\o\do\p\do\q\do\r\do\s\do\t%
  \do\u\do\v\do\w\do\x\do\y\do\z\do\A\do\B\do\C\do\D%
  \do\E\do\F\do\G\do\H\do\I\do\J\do\K\do\L\do\M\do\N%
  \do\O\do\P\do\Q\do\R\do\S\do\T\do\U\do\V\do\W\do\X%
  \do\Y\do\Z}%
\fi
%    \end{macrocode}
%
%
%
%
% \subsection{Metainformationen}
%

%
% \subsubsection{Autor}
%
% \begin{macro}{\authorwithand}
%   Im Autor-Makro werden mehrere Autoren durch das Makro \verb|\and|
%   getrennt. Mit dem Makro \verb|\authorwithand| können die Autoren
%   komfortabel mit einem beliebigen Trennzeichen (Wert von
%   \verb|\and|) gedruckt werden. Das Makro unterdrückt die Angaben in
%   \verb|\thanks|.
%
%    \begin{macrocode}
\newcommand*\authorwithand[1]{%
  \let\tempand=\and%
  \renewcommand*\and{\leavevmode\unskip#1}%
  \let\tempthanks=\thanks%
  \renewcommand*\thanks[1]{}%
  \@author%
  \let\and=\tempand%
  \let\thanks=\tempthanks%
}
%    \end{macrocode}
%
% \end{macro}
%
% \subsubsection{Schlüsselwörter}
%
% \begin{macro}{\keywords}
% \begin{macro}{\@keywords}
%   Zusätzlich zu den üblichen \LaTeX-Metadaten bieten unsere Klassen
%   auch Zugriff auf das Schlüsselwort-Feld der
%   PDF-Dokumenteninformationen. Dem Benutzer wird dazu das Makro
%   \verb|\keywords| zur Verfügung gestellt. Wird das hyperref-Paket
%   geladen, dann sollen die Schlüsselwörter auch automatisch in den
%   PDF-Dokumenteneigenschaften stehen; andernfalls sind die
%   Schlüsselwörter nur zum Abdrucken im Rahmen der Titelei gedacht.
%
%    \begin{macrocode}
\newcommand*{\@keywords}{}
\ifldv@doloadhyperref%
\newcommand{\keywords}[1]{%
  \renewcommand*\@keywords{#1}%
  \hypersetup{pdfkeywords = {#1}}%
}%
\else%
\newcommand{\keywords}[1]{%
  \renewcommand*\@keywords{#1}%
}%
\fi%
%    \end{macrocode}
% \end{macro}
% \end{macro}
%
% \begin{macro}{\keywordsname}
%   Um die Schlüsselwörter zu setzen, brauchen wir noch den
%   sprachabhängigen Begriff für \enquote{Schlüsselwörter}. Die unten
%   gezeigte Vorgehensweise integriert sich in den Mechanismus von
%   Babel; der Vorgabetext ist der englische.
%
%    \begin{macrocode}
\newcommand*\keywordsname{Key words}%
\addto\captionsngerman{%
  \renewcommand*\keywordsname{Schlagw\"orter}%
}
\addto\captionsgerman{%
  \renewcommand*\keywordsname{Schlagw\"orter}%
}
\addto\captionsenglish{%
  \renewcommand*\keywordsname{Key words}%
}
%    \end{macrocode}
% \end{macro}
%
% \begin{macro}{\makekeywords}
%   Im Text können die Schlüsselwörter dann mit dem Befehl
%   \verb|\makekeywords| eingefügt werden. Diese bei Artikeln
%   geschieht üblicherweise direkt nach dem Abstract, bei größeren
%   Dokumenten in der Nähe der Titelseite bei anderen Metadaten.
%
%    \begin{macrocode}
\newcommand{\makekeywords}{%
  \if@twocolumn
  \vspace{\topsep}
  \noindent{\sectfont\size@paragraph\keywordsname:} \@keywords
  \else
  \small
  \begin{center}%
    {\normalfont\sectfont\nobreak\keywordsname
      \vspace{-.5em}\vspace{\z@}}%
  \end{center}%
  \quotation\@keywords\endquotation
  \fi}
%    \end{macrocode}
% \end{macro}
%
%
% \subsubsection{Dokumentenversion}
%
% \begin{macro}{\version}
% \begin{macro}{\@version}
%   Neben der Datumsinformation wird Dokumenten häufig auch eine
%   Versionsnummer zugewiesen, um dessen Fortentwicklung zu
%   kennzeichnen. Der Benutzer kann diese mit dem Makro
%   \verb|\version| einstellen; sie erscheint dann auf der Titelseite
%   und in der Fußzeile.
%
%    \begin{macrocode}
\newcommand*{\@version}{}
\newcommand{\version}[1]{\renewcommand*\@version{#1}}
%    \end{macrocode}
% \end{macro}
% \end{macro}
%
%
% \subsubsection{Verlag: Universität und Lehrstuhl}
%
% Die \KOMAScript-Klassen definieren das Verlagsfeld durch das Makro
% \verb|\publishers|. Bei uns ist es gleich sinnvoll vorbelegt und
% wird an verschiedenen Stellen verwendet.
%
%    \begin{macrocode}
\publishers{Technische Universit\"at M\"unchen}
%    \end{macrocode}
%
% \begin{macro}{\institute}
% \begin{macro}{\@institute}
% \begin{macro}{\postadress}
% \begin{macro}{\@postadress}
% \begin{macro}{\citationaddress}
% \begin{macro}{\@citationaddress}
% \begin{macro}{\publishersurl}
% \begin{macro}{\@publishersurl}
%   Zusätzlich sind auch der Lehrstuhlname, die Postadresse, der Ort
%   für den Referenztext und die URL des Verlags/Lehrstuhls
%   einstellbar:
%
%    \begin{macrocode}
\newcommand*\@institute{Chair of Data Processing}
\newcommand*\institute[1]{\renewcommand*\@institute{#1}}
\newcommand*{\@postaddress}{80290 M\"unchen, Germany}
\newcommand*{\postaddress}[1]{\renewcommand*{\@postaddress}{#1}}
\newcommand*{\@citationaddress}{Munich, Germany}
\newcommand*{\citationaddress}[1]{\renewcommand*{\@citationaddress}{#1}}
\newcommand*{\@publishersurl}{\url{http://www.ldv.ei.tum.de/}}
\newcommand*{\publishersurl}[1]{\renewcommand*{\@publishersurl}{#1}}
\addto\captionsenglish{%
  \institute{Chair of Data Processing}
  \postaddress{80290 M\"unchen, Germany}
  \citationaddress{Munich, Germany}
}
\addto\captionsngerman{%
  \institute{Lehrstuhl f\"ur Datenverarbeitung}
  \postaddress{80290 M\"unchen}
  \citationaddress{M\"unchen}
}
\addto\captionsgerman{%
  \institute{Lehrstuhl f\"ur Datenverarbeitung}
  \postaddress{80290 M\"unchen}
  \citationaddress{M\"unchen}
}
%    \end{macrocode}
% \end{macro}
% \end{macro}
% \end{macro}
% \end{macro}
% \end{macro}
% \end{macro}
% \end{macro}
% \end{macro}
%
%
% \subsubsection{Betreuer einer studentischen Abschlussarbeit}
%
% \begin{macro}{\supervisor}
% \begin{macro}{\@supervisor}
%   Für ein Diplom-, Master-, Studien- oder Bachelorarbeit kann der
%   betreuende Assistent angegeben werden. Diese Information wird dann
%   in der Titelei abgedruckt.
%
% \begin{macrocode}
%<*book>
\ifldv@studthesis
  \newcommand*\@supervisor{}
  \newcommand*\supervisor[1]{\renewcommand*\@supervisor{#1}}
\fi
%</book>
%    \end{macrocode}
% \end{macro}
% \end{macro}
%
%
% \subsubsection{Lizenz}
%
% \begin{macro}{\licensetext}
% \begin{macro}{\ldv@licensetext}
%   Der Benutzer kann das Werk unter eine Lizenz stellen. Dieser Text
%   wird auf der Impressumsseite oder, bei einseitigem Druck, auf der
%   Titelseite abgedruckt.  Der Text kann vom Benutzer individuell mit
%   dem Makro \verb|licensetext| bestimmt werden.  Er wird intern in
%   \verb|ldv@licensetext| gespeichert.
%
%    \begin{macrocode}
\newcommand*\ldv@licensetext{}
\newcommand\licensetext[1]{\renewcommand*\ldv@licensetext{#1}}
%    \end{macrocode}
% \end{macro}
% \end{macro}
%
% \begin{macro}{\license}
%   Mit dem Makro \verb|\license| kann der Benutzer einfach eine der
%   hinterlegten Lizenzen auswählen, ohne sich um den genauen
%   Lizenztext kümmern zu müssen. Im Moment sind die Text der sechs
%   Creative Commons-Lizenzen für Deutschland verfügbar: CC-BY,
%   CC-BY-SA, CC-BY-ND, CC-BY-NC, CC-BY-NC-SA, CC-BY-NC-ND. Je nach
%   ausgewählter Lizenz und aktiver Sprache wird der passende Text in
%   \verb|\ldv@licensetext| abgelegt.
%
%   Zur einfachen Implementierung greift dieses Makro auf einen choice
%   key des Pakets xkeyval zurück. Es ist der choice key
%   \emph{license}, welcher unten implementiert ist.
%
%    \begin{macrocode}
\newcommand*\license[1]{\setkeys{ldv}{license=#1}}
%    \end{macrocode}
% \end{macro}
%
% \begin{option}{license}
% \begin{macro}{\ldv@licensetext@ccby}
% \begin{macro}{\ldv@licensetext@ccbysa}
% \begin{macro}{\ldv@licensetext@ccbynd}
% \begin{macro}{\ldv@licensetext@ccbync}
% \begin{macro}{\ldv@licensetext@ccbyncsa}
% \begin{macro}{\ldv@licensetext@ccbyncnd}
%   Mit dem choice key \emph{license} lässt sich das Makro
%   \verb|\license| sehr einfach implementieren.  Dazu speicheren die
%   Makros \verb|\ldv@license@...| den jeweiligen Lizenztext in der
%   aktuellen Sprache.  Bei der Verarbeitung der Option license wird
%   dann nur noch \verb|\ldv@licensetext| auf das jeweilige Makro der
%   Lizenz gelinkt.
%
%    \begin{macrocode}
\newcommand*\ldv@licensetext@ccby{This work is licensed under the Creative Commons Attribution 4.0 International License. To view a copy of this license, visit \url{http://creativecommons.org/licenses/by/4.0/} or send a letter to Creative Commons, PO Box 1866, Mountain View, CA 94042, USA.}%
\newcommand*\ldv@licensetext@ccbysa{This work is licensed under the Creative Commons Attribution-ShareAlike 4.0 International License. To view a copy of this license, visit \url{http://creativecommons.org/licenses/by-sa/4.0/} or send a letter to Creative Commons, PO Box 1866, Mountain View, CA 94042, USA.}%
\newcommand*\ldv@licensetext@ccbynd{This work is licensed under the Creative Commons Attribution-NoDerivatives 4.0 International License. To view a copy of this license, visit \url{http://creativecommons.org/licenses/by-nd/4.0/} or send a letter to Creative Commons, PO Box 1866, Mountain View, CA 94042, USA.}%
\newcommand*\ldv@licensetext@ccbync{This work is licensed under the Creative Commons Attribution-NonCommercial 4.0 International License. To view a copy of this license, visit \url{http://creativecommons.org/licenses/by-nc/4.0/} or send a letter to Creative Commons, PO Box 1866, Mountain View, CA 94042, USA.}%
\newcommand*\ldv@licensetext@ccbyncsa{This work is licensed under the Creative Commons Attribution-NonCommercial-ShareAlike 4.0 International License. To view a copy of this license, visit \url{http://creativecommons.org/licenses/by-nc-sa/4.0/} or send a letter to Creative Commons, PO Box 1866, Mountain View, CA 94042, USA.}%
\newcommand*\ldv@licensetext@ccbyncnd{This work is licensed under the Creative Commons Attribution-NonCommercial-NoDerivatives 4.0 International License. To view a copy of this license, visit \url{http://creativecommons.org/licenses/by-nc-nd/4.0/} or send a letter to Creative Commons, PO Box 1866, Mountain View, CA 94042, USA.}%
\addto\captionsenglish{%
  \renewcommand*\ldv@licensetext@ccby{This work is licensed under the Creative Commons Attribution 4.0 International License. To view a copy of this license, visit \url{http://creativecommons.org/licenses/by/4.0/} or send a letter to Creative Commons, PO Box 1866, Mountain View, CA 94042, USA.}%
  \renewcommand*\ldv@licensetext@ccbysa{This work is licensed under the Creative Commons Attribution-ShareAlike 4.0 International License. To view a copy of this license, visit \url{http://creativecommons.org/licenses/by-sa/4.0/} or send a letter to Creative Commons, PO Box 1866, Mountain View, CA 94042, USA.}%
  \renewcommand*\ldv@licensetext@ccbynd{This work is licensed under the Creative Commons Attribution-NoDerivatives 4.0 International License. To view a copy of this license, visit \url{http://creativecommons.org/licenses/by-nd/4.0/} or send a letter to Creative Commons, PO Box 1866, Mountain View, CA 94042, USA.}%
  \renewcommand*\ldv@licensetext@ccbync{This work is licensed under the Creative Commons Attribution-NonCommercial 4.0 International License. To view a copy of this license, visit \url{http://creativecommons.org/licenses/by-nc/4.0/} or send a letter to Creative Commons, PO Box 1866, Mountain View, CA 94042, USA.}%
  \renewcommand*\ldv@licensetext@ccbyncsa{This work is licensed under the Creative Commons Attribution-NonCommercial-ShareAlike 4.0 International License. To view a copy of this license, visit \url{http://creativecommons.org/licenses/by-nc-sa/4.0/} or send a letter to Creative Commons, PO Box 1866, Mountain View, CA 94042, USA.}%
  \renewcommand*\ldv@licensetext@ccbyncnd{This work is licensed under the Creative Commons Attribution-NonCommercial-NoDerivatives 4.0 International License. To view a copy of this license, visit \url{http://creativecommons.org/licenses/by-nc-nd/4.0/} or send a letter to Creative Commons, PO Box 1866, Mountain View, CA 94042, USA.}%
}
\addto\captionsngerman{%
  \renewcommand*\ldv@licensetext@ccby{Dieses Werk ist unter einer Creative Commons Lizenz vom Typ Namensnennung 4.0 International zugänglich. Um eine Kopie dieser Lizenz einzusehen, konsultieren Sie \url{http://creativecommons.org/licenses/by/4.0/} oder wenden Sie sich brieflich an Creative Commons, Postfach 1866, Mountain View, California, 94042, USA.}%
  \renewcommand*\ldv@licensetext@ccbysa{Dieses Werk ist unter einer Creative Commons Lizenz vom Typ Namensnennung - Weitergabe unter gleichen Bedingungen 4.0 International zugänglich. Um eine Kopie dieser Lizenz einzusehen, konsultieren Sie \url{http://creativecommons.org/licenses/by-sa/4.0/} oder wenden Sie sich brieflich an Creative Commons, Postfach 1866, Mountain View, California, 94042, USA.}%
  \renewcommand*\ldv@licensetext@ccbynd{Dieses Werk ist unter einer Creative Commons Lizenz vom Typ Namensnennung - Keine Bearbeitungen 4.0 International zugänglich. Um eine Kopie dieser Lizenz einzusehen, konsultieren Sie \url{http://creativecommons.org/licenses/by-nd/4.0/} oder wenden Sie sich brieflich an Creative Commons, Postfach 1866, Mountain View, California, 94042, USA.}%
  \renewcommand*\ldv@licensetext@ccbync{Dieses Werk ist unter einer Creative Commons Lizenz vom Typ Namensnennung-Nicht kommerziell 4.0 International zugänglich. Um eine Kopie dieser Lizenz einzusehen, konsultieren Sie \url{http://creativecommons.org/licenses/by-nc/4.0/} oder wenden Sie sich brieflich an Creative Commons, Postfach 1866, Mountain View, California, 94042, USA.}%
  \renewcommand*\ldv@licensetext@ccbyncsa{Dieses Werk ist unter einer Creative Commons Lizenz vom Typ Namensnennung - Nicht-kommerziell - Weitergabe unter gleichen Bedingungen 4.0 International zugänglich. Um eine Kopie dieser Lizenz einzusehen, konsultieren Sie \url{http://creativecommons.org/licenses/by-nc-sa/4.0/} oder wenden Sie sich brieflich an Creative Commons, Postfach 1866, Mountain View, California, 94042, USA.}%
  \renewcommand*\ldv@licensetext@ccbyncnd{Dieses Werk ist unter einer Creative Commons Lizenz vom Typ Namensnennung - Nicht kommerziell - Keine Bearbeitungen 4.0 International zugänglich. Um eine Kopie dieser Lizenz einzusehen, konsultieren Sie \url{http://creativecommons.org/licenses/by-nc-nd/4.0/} oder wenden Sie sich brieflich an Creative Commons, Postfach 1866, Mountain View, California, 94042, USA.}%
}
\define@choicekey*+{ldv}{license}[\ldv@choicekeyval\ldv@choicekeynr]{cc-by,cc-by-sa,cc-by-nd,cc-by-nc,cc-by-nc-sa,cc-by-nc-nd}{%
  \ifcase\ldv@choicekeynr\relax
    % Namensnennung: CC-BY
    \renewcommand*{\ldv@licensetext}{\ldv@licensetext@ccby}
  \or
    % Namensnennung und Weitergabe unter gleichen Bedingungen: CC-BY-SA
    \renewcommand*{\ldv@licensetext}{\ldv@licensetext@ccbysa}
  \or
    % Namensnennung und keine Bearbeitung: CC-BY-ND
    \renewcommand*{\ldv@licensetext}{\ldv@licensetext@ccbynd}
  \or
    % Namensnennung und nicht-kommerzielle Nutzung: CC-BY-NC
    \renewcommand*{\ldv@licensetext}{\ldv@licensetext@ccbync}
  \or
    % Namensnennung, nicht-kommerzielle Nutzung und Weitergabe unter
    % gleichen Bedingungen: CC-BY-NC-SA
    \renewcommand*{\ldv@licensetext}{\ldv@licensetext@ccbyncsa}
  \or
    % Namensnennung, nicht-kommerzielle Nutzung und keine Bearbeitung:
    % CC-BY-NC-ND
    \renewcommand*{\ldv@licensetext}{\ldv@licensetext@ccbyncnd}
  \fi
}{%
  \ClassWarning{\@currname}{%
    Value given for the option license is not known.%
  }
}
%    \end{macrocode}
% \end{macro}
% \end{macro}
% \end{macro}
% \end{macro}
% \end{macro}
% \end{macro}
% \end{option}
%
%
%
% \subsubsection{Dokumenteninformationen in PDF-Dateien}
%
% \begin{macro}{\title}
% \begin{macro}{\subtitle}
% \begin{macro}{\author}
% \begin{macro}{\subject}
%   Die Metainformationen sollen automatisch in die
%   Dokumenteneigenschaften der PDF-Datei eingetragen werden. Dies
%   geschieht am Besten zu dem Zeitpunkt, zu dem der Autor die
%   jeweilge Metainformation setzt. Das neu eingeführte Metaelement
%   \enquote{keywords} tut das schon. Für alle bestehenden Metaelemente
%   müssen die Befehle dazu noch umdefiniert werden. Dies passiert
%   natürlich nur, wenn das hyperref-Paket tatsächlich geladen wird.
%
%    \begin{macrocode}
\ifldv@doloadhyperref%
\renewcommand*\title[1]{%
  \renewcommand*\@title{#1}%
  \hypersetup{%
    pdftitle = {\@title\ifx\@subtitle\@empty\else. \@subtitle.\fi}%
  }%
}
\renewcommand*\subtitle[1]{%
  \renewcommand*\@subtitle{#1}%
  \hypersetup{%
    pdftitle = {\ifx\@title\@empty\else\@title. \@subtitle.\fi}%
  }%
}
\renewcommand*\author[1]{%
  \renewcommand*\@author{#1}%
  \let\tempand=\and%
  \renewcommand*\and{, }%
  \let\tempthanks=\thanks%
  \let\thanks=\@empty%
  \hypersetup{%
    pdfauthor  = {\@author}%
  }%
  \let\and=\tempand%
  \let\thanks=\tempthanks%
}
\renewcommand*\subject[1]{%
  \renewcommand*\@subject{#1}%
  \hypersetup{%
    pdfsubject = {\@subject}
  }%
}
\fi%
%    \end{macrocode}
% \end{macro}
% \end{macro}
% \end{macro}
% \end{macro}
%
%
%
% \subsection{Titelei}
%
% \subsubsection{Grundeinstellungen}
%
% Zu Anfang wird die \KOMAScript-Schrift für das Subject-Feld
% angepasst, so dass es jetzt in normaler Schriftgröße erscheint. Dies
% wirkt sich auf alle Titel-Makros aus.
%
%    \begin{macrocode}
\renewcommand*\subject@font{\normalfont\normalcolor\bfseries}
%    \end{macrocode}
%
% \begin{macro}{\ldv@ldvlogowidth}
% \begin{macro}{\ldv@ldvlogoS}
% \begin{macro}{\ldv@ldvlogoCMYK}
% \begin{macro}{\ldv@tumlogowidth}
% \begin{macro}{\ldv@tumlogoS}
% \begin{macro}{\ldv@tumlogoCMYK}
%   Dateinamen und die Breiten der Logos werden auch von mehreren
%   Makros im Folgenden verwendet werden. Bisher sind es lediglich die
%   Breiten für ein A4-Papier. Das TUM-Logo ist bei dieser Breite
%   knapp 10mm hoch.
%
%    \begin{macrocode}
\RequirePackage{calc}
\newlength\ldv@ldvlogowidth
\setlength\ldv@ldvlogowidth{19mm}
\newcommand*\ldv@ldvlogoS{LDVLogoS_oT}
\newcommand*\ldv@ldvlogoCMYK{LDVLogoCMYK_oT}
\newlength\ldv@tumlogowidth
\setlength\ldv@tumlogowidth{17mm}
\newcommand*\ldv@tumlogoS{TUMLogo_oZ_Vollfl_sw}
\newcommand*\ldv@tumlogoCMYK{TUMLogo_oZ_Vollfl_CMYK}
%    \end{macrocode}
% \end{macro}
% \end{macro}
% \end{macro}
% \end{macro}
% \end{macro}
% \end{macro}
%
% \begin{macro}{\andname}
%   Zur Lokalisierung des Autor-Feldes wird der sprachspezifische
%   Begriff für \enquote{und} benötigt.
%    \begin{macrocode}
%<*book>
  \newcommand*\andname{and}
  \addto\captionsngerman{%
    \renewcommand*\andname{und}%
  }
  \addto\captionsgerman{%
    \renewcommand*\andname{und}%
  }
  \addto\captionsenglish{%
    \renewcommand*\andname{and}%
  }
%</book>
%    \end{macrocode}
% \end{macro}
%
% \begin{macro}{\diplomarbeitname}
% \begin{macro}{\mastersthesisname}
% \begin{macro}{\studienarbeitname}
% \begin{macro}{\bachelorsthesisname}
% \begin{macro}{\idpname}
% \begin{macro}{\phdthesisname}
%   Für studentische Abschlussarbeiten und Dissertationen soll das
%   \verb|subject| den genauen Typ der Arbeit enthalten. Dieser kann
%   eine Diplom-, Master-, Studien- oder Bachelorarbeit, ein
%   interdisziplinäres Projekt sowie eine Disseration sein. Die in
%   Kapitel~\ref{sec:deklaration-optionen} definierte Bedingung
%   \verb|\ifldv@studthesis| trifft zu, falls in den Klassenoptionen
%   eine dieser studentischen Abschlussarbeiten ausgewählt wurde.
%   \verb|\ifldv@phdthesis| trifft zu, falls in den Klassenoptionen
%   eine Disseration ausgewählt wurde.
%
%    \begin{macrocode}
%<*book>
\ifldv@studthesis%
  \newcommand*\diplomarbeitname{Diploma thesis}
  \newcommand*\mastersthesisname{Master's thesis}
    \newcommand*\studienarbeitname{Thechnical report}
  \newcommand*\bachelorsthesisname{Bachelor's thesis}
  \newcommand*\idpname{Interdisciplinary project}
  \addto\captionsngerman{%
    \renewcommand*\diplomarbeitname{Diplomarbeit}%
    \renewcommand*\mastersthesisname{Masterarbeit}%
    \renewcommand*\studienarbeitname{Studienarbeit}%
    \renewcommand*\bachelorsthesisname{Bachelorarbeit}%
    \renewcommand*\idpname{Interdisziplin\"ares Projekt}%
  }
  \addto\captionsgerman{%
    \renewcommand*\diplomarbeitname{Diplomarbeit}%
    \renewcommand*\mastersthesisname{Masterarbeit}%
    \renewcommand*\studienarbeitname{Studienarbeit}%
    \renewcommand*\bachelorsthesisname{Bachelorarbeit}%
    \renewcommand*\idpname{Interdisziplin\"ares Projekt}%
  }
  \addto\captionsenglish{%
    \renewcommand*\diplomarbeitname{Diploma thesis}%
    \renewcommand*\mastersthesisname{Master's thesis}%
    \renewcommand*\studienarbeitname{Technical repor}%
    \renewcommand*\bachelorsthesisname{Bachelor's thesis}%
    \renewcommand*\idpname{Interdisciplinary project}%
  }
  \ifstr{\ldv@doctype}{diplomarbeit}{
    \subject{\diplomarbeitname}
  }{\ifstr{\ldv@doctype}{mastersthesis}{
      \subject{\mastersthesisname}
    }{\ifstr{\ldv@doctype}{idp}{
        \subject{\idpname}
      }{\ifstr{\ldv@doctype}{bachelorsthesis}{
          \subject{\bachelorsthesisname}
        }{
          \subject{\studienarbeitname}
        }
      }
    }
  }
\fi
\ifldv@phdthesis
  \newcommand*\phdthesisname{Dissertation}
  \addto\captionsenglish{%
    \renewcommand*\phdthesisname{Dissertation}
  }
  \addto\captionsngerman{%
    \renewcommand*\phdthesisname{Dissertation}
  }
  \addto\captionsngerman{%
    \renewcommand*\phdthesisname{Dissertation}
  }
  \subject{\phdthesisname}
\fi
%</book>
%    \end{macrocode}
% \end{macro}
% \end{macro}
% \end{macro}
% \end{macro}
% \end{macro}
% \end{macro}
%
%
% \subsubsection{Setzen der Titelei}
%
% Das Makro \verb|\maketitle| setzt die gesamte Titelei; als Basis
% habe ich die Definition aus den \KOMAScript-Quellen genommen und
% weiter modularisiert. Der Titel selbst (insbesondere die Titelseite)
% wird jetzt von eigenen Makros implementiert, so dass man sich beim
% Redefinieren rein auf das Layout konzentrieren kann, ohne die
% umliegenden Seiten beachten zu müssen. So können Benutzern leichter
% einen eigenen Titel definieren.
%
% \verb|\maketitle| benutzt nun den Schlüssel-Wert-Mechanismus des
% \emph{xkeyval}-Pakets, um flexibel vielfältige Funktionen
% realisieren zu können. Das Makro unterstützt folgende Optionen:
% \begin{itemize}
% \item \verb|frontcover| (z.B. \verb|frontcover=Design1|)
% \item \verb|pagenumber| (z.B. \verb|pagenumber=3|)
% \end{itemize}
% Zuerst werden im Folgenden die Optionen definiert.
%
% \begin{option}{frontcover}
% \begin{macro}{\ldv@frontcoverdesign}
%   Die Option \verb|frontcover| gibt an, welches Design aus einer
%   vordefinierten Auswahl für den vorderen Umschlag verwendet werden
%   soll. Die verschiedenen Umschlagseiten sollen ohne viel Aufwand
%   für den Autor ein gewisses Maß an Individualität bieten. Der
%   Benutzer wählt ein Design in der Form \verb|frontcover=Design1|
%   aus. Ohne diese Option erscheint keine Umschlagseite; damit ist
%   \verb|maketitle| weiterhin kompatibel zu den
%   Standard-\LaTeX-Klassen.
%
%    \begin{macrocode}
\define@choicekey*{ldv}{frontcover}[\ldv@frontcoverdesign]
  {design1}{%
    \if@titlepage\else%
      \ClassWarning{\@currname}{%
        Option frontcover of \string\maketitle\ is only valid
        \MessageBreak
        when using titlepage=true%
      }%
    \fi%
  }%
%    \end{macrocode}
% \end{macro}
% \end{option}
%
% \begin{option}{frontcoverfile}
%
%    \begin{macrocode}
\define@cmdkey{ldv}[ldv@]{frontcoverfile}{}
%    \end{macrocode}
% \end{option}
%
% \begin{option}{pagenumber}
%   Die Option \verb|pagenumber| ermöglicht es dem Autor, eine
%   bestimmte Seitenzahl für den Beginn der Titelei anzugeben. Dies
%   ist eine Option, die das \KOMAScript-Makro als optionales Argument
%   bietet und von dort übernommen wurde, um etwa funktionsgleich zu
%   bleiben. Ohne diese Option beginnt die Zählung bei 1.
%
%    \begin{macrocode}
\define@key{ldv}{pagenumber}{%
  \if@titlepage%
    \setcounter{page}{#1}%
  \else%
        \ClassWarning{\@currname}{%
      Option pagenumber of \string\maketitle\ is only valid
      \MessageBreak
      when using titlepage=true%
    }%
  \fi%
}
%    \end{macrocode}
% \end{option}
%
% \begin{macro}{\@maketitle}
%   Das Makro \verb|\@maketitle| aus den \KOMAScript-Klassen benutze
%   ich nicht mehr. Deshalb deaktiviere ich es hier.
%
%    \begin{macrocode}
\global\let\@maketitle\relax
%    \end{macrocode}
% \end{macro}
%
% \begin{macro}{\maketitle}
%   Wie in den \KOMAScript-Klassen auch gibt es jeweils eine eigene
%   Makro-Definition für die Option einer ganzen Titelseite und eines
%   einfachen Titelkopfes. Die komplette Titelei steht nur für die
%   Option \verb|titlepage=true| zur Verfügung.
%
%   Das Makro nutzt in diesen LDV-Klassen den
%   Schlüssel-Wert-Mechanismus des \emph{xkeyval}-Pakets. Die Optionen
%   sind weiter oben in diesem Kapitel beschrieben.
%
%   Schließlich löscht dieses Makro am Ende nun nicht mehr alle
%   Makros zu den Meta-Informationen. Diese kann der Autor also bei
%   diesen Klassen im Nachhinein verwenden.
%
%    \begin{macrocode}
\if@titlepage
  \renewcommand*\maketitle[1][]{%
    \setkeys{ldv}{#1}%
    \begin{titlepage}
      \let\footnotesize\small
      \let\footnoterule\relax
      \let\footnote\thanks
      \renewcommand*\thefootnote{\@fnsymbol\c@footnote}%
      \let\@oldmakefnmark\@makefnmark
      \renewcommand*{\@makefnmark}{\rlap\@oldmakefnmark}%
      \ifdefined\ldv@frontcoverfile%
        \setcounter{page}{-1}%
        \ldv@includefile{\ldv@frontcoverfile}%
        \if@twoside\cleardoubleemptypage\else\clearpage\fi%
        \thispagestyle{empty}
      \else%
        \ifdefined\ldv@frontcoverdesign%
          \ifstr{\ldv@frontcoverdesign}{design1}{%
            \setcounter{page}{-1}%
            \ldv@makecover@eins%
            \if@twoside\cleardoubleemptypage\else\clearpage\fi%
            \thispagestyle{empty}
          }{}%
        \fi%
      \fi%
      \ifx\@extratitle\@empty \else
        \noindent\@extratitle\next@tpage\cleardoubleemptypage
        \thispagestyle{empty}%
      \fi
%<+book>      \ifldv@phdthesis\ldv@phdtitle\else\ldv@pagetitle\fi%
%<+article>      \ldv@pagetitle%
      \if@twoside\next@tpage
        \begin{minipage}[t]{\textwidth}
          \@uppertitleback
        \end{minipage}\par
        \vfill
        \begin{minipage}[b]{\textwidth}
          \@lowertitleback
        \end{minipage}
      \fi
      \ifx\@dedication\@empty \else
        \next@tpage\null\vfill
        {\centering \Large \@dedication \par}%
        \vskip \z@ \@plus3fill
        \if@twoside \next@tpage\cleardoubleemptypage \fi
      \fi
    \end{titlepage}
    \setcounter{footnote}{0}%
    \global\let\thanks\relax
    \global\let\maketitle\relax
    \global\let\@thanks\@empty
  }
\else
  \renewcommand*\maketitle[1][]{\par%
    \setkeys{ldv}{#1}%
    \begingroup
      \renewcommand*\thefootnote{\@fnsymbol\c@footnote}%
      \let\@oldmakefnmark\@makefnmark
      \renewcommand*{\@makefnmark}{\rlap\@oldmakefnmark}
      \if@twocolumn
        \ifnum \col@number=\@ne
          \ldv@headtitle
         \else
          \twocolumn[\ldv@headtitle]%
        \fi
      \else
        \newpage
        \global\@topnum\z@
        \ldv@headtitle
      \fi
      \thispagestyle{\titlepagestyle}\@thanks
    \endgroup
    \setcounter{footnote}{0}%
    \let\thanks\relax
    \let\maketitle\relax
    \global\let\@thanks\@empty
  }
\fi
%    \end{macrocode}
% \end{macro}
%
%
% \subsubsection{Titelseite}
%
% \begin{macro}{\@titletitle}
% \begin{macro}{\titletitle}
%   Um für die Titelseite eigene Zeilenumbrüche im Titel zu setzen,
%   kann der Benutzer optional das Hilfsmakro \verb|\titletitle|
%   benutzen. Der so festgelegte Titel wird dann beim Setzen der
%   Titelseite bevorzugt.
%
%    \begin{macrocode}
\newcommand*\@titletitle{}
\newcommand*\titletitle[1]{%
  \renewcommand*\@titletitle{#1}%
}
%    \end{macrocode}
% \end{macro}
% \end{macro}
%
% \begin{macro}{\ldv@studthesispersons}
%   Bei studentischen Abschlussarbeiten mit einseitigem Druck wird der
%   Betreuer auf der Titelseite genannt. Dazu ist ein zusätzliches
%   lokalisiertes Makro nötig, welches anstelle der Autoren gedruckt
%   werden soll. (Die Implementierung hier kann mit mehreren Autoren
%   umgehen, obwohl das bei einer Abschlussarbeit mit derzeitiger
%   Prüfungsordnung wenig Sinn macht.)
%
%    \begin{macrocode}
%<*book>
\ifldv@studthesis
  \if@twoside\else
    \newcommand*\ldv@studthesispersons{%
      Written by \authorwithand{\andname}\\
      as a project final report of the course  Prof.\ Dr.-Ing.\ Klaus Diepold%
      \ifx\@supervisor\@empty\else \ and \@supervisor\fi%
    }
    \addto\extrasenglish{
      \renewcommand*\ldv@studthesispersons{%
        Written by \authorwithand{\andname}\\
        Supervised by Prof.\ Dr.-Ing.\ Klaus Diepold%
        \ifx\@supervisor\@empty\else \ and \@supervisor\fi%
      }
    }
    \addto\extrasngerman{%
      \renewcommand*\ldv@studthesispersons{%
        Verfasst von \authorwithand{\andname}\\
        Betreut von Prof.\ Dr.-Ing.\ Klaus Diepold%
        \ifx\@supervisor\@empty\else \ und \@supervisor\fi%
      }%
    }
    \addto\extrasgerman{%
      \renewcommand*\ldv@studthesispersons{%
        Verfasst von \authorwithand{\andname}\\
        Betreut von Prof.\ Dr.-Ing.\ Klaus Diepold%
        \ifx\@supervisor\@empty\else \ und \@supervisor\fi%
        \let\and=\tempand%
      }%
    }
  \fi
\fi
%</book>
%    \end{macrocode}
% \end{macro}
%
% \begin{macro}{\ldv@pagetile}
%   Mit dem Makro \verb|ldv@pagetitle| wird ein Titel gesetzt, der
%   eine ganze Seite in Anspruch nimmt. Der Code wurde im Wesentlichen
%   aus den \KOMAScript-Quellen entnommen (aus dem Makro
%   \verb|\maketitle|). Folgende Änderungen sind dann eingeflossen:
%
%   \begin{itemize}
%   \item Coporate Design-Elemente eingefügt (Logo und Namen).
%   \item Abstand zwischen Untertitel und Autoren korrigiert.
%   \item Verlags-/Institutions-Angabe entfernt; sie ist bereits durch die oben
%     genannten Coporate Design-Elemente vorhanden.
%   \item Versionsnummer eingefügt.
%   \end{itemize}
%
%    \begin{macrocode}
\if@titlepage
\newcommand*\ldv@pagetitle{%
  \ldv@makepublishers%
  \setparsizes{\z@}{\z@}{\z@\@plus 1fil}\par@updaterelative%
  \ifx\@titlehead\@empty \else%
  \begin{minipage}[t]{\textwidth}%
    \@titlehead
  \end{minipage}\par
  \fi
  \null\vfill
  \begin{center}
    \ifx\@subject\@empty\else
    {\subject@font \@subject \par}%
    \vskip 3em
    \fi
    {\titlefont\huge \ifx\@titletitle\@empty \@title \else\@titletitle \fi\par}%
    \ifx\@subtitle\@empty\else%
      \vskip 1em%
      {\usekomafont{subtitle}\@subtitle\par}\fi%
    \vskip 3em
    {\lineskip 0.75em%
%<*book>
      \ifdefined\ldv@studthesispersons%
      \ldv@studthesispersons\par%
      \else%
%</book>
      \begin{tabular}[t]{c}
        \@author
      \end{tabular}\par%
%<+book>      \fi
    }%
    \vskip 1.5em
    {\@date%
      \ifx\@version\@empty\else, \@version\fi\par}%
    {\ifx\@license\@empty\else\if@twoside\else%
      \vskip 1.5em%
      \ldv@licensetext\par%
    \fi\fi}%
    \vfill\vfill\vfill\vfill\null
  \end{center}\par
  \@thanks
}%
%    \end{macrocode}
% \end{macro}
%
% \begin{macro}{\ldv@makepublishers}
%   \verb|\ldv@makepublishers| zeichnet am Fuß der Seite die
%   Lehrstuhl-In\-for\-ma\-tio\-nen (Logo und Name). Dieser Vorgang
%   ist aufwändiger, so dass er in einem eigenen Makro implementiert
%   ist.
%
%   Zuerst wird der Inhalt in einer Box (\verb|\ldv@pubslishersbox|)
%   angeordnet. Diese wird dann an den Fuß des Blatts
%   verschoben. Dabei soll der untere Rand (\verb|ldv@bottommargin|)
%   genauso groß sein wie der aktuell eingestellt linke Rand. Die
%   Verschiebung ist in der Länge \verb|\ldv@publishersmove|
%   gespeichert.
%    \begin{macrocode}
\newcommand*\ldv@makepublishers{%
  \newlength\ldv@tumwidth%
  \newlength\ldv@ldvwidth%
  \newlength\ldv@publisherwidth%
  \settowidth\ldv@tumwidth{\@publishers}%
  \settowidth\ldv@ldvwidth{\@institute}%
  \ifdim\ldv@tumwidth >\ldv@ldvwidth%
    \setlength\ldv@publisherwidth{\ldv@tumwidth}%
  \else%
    \setlength\ldv@publisherwidth{\ldv@ldvwidth}%
  \fi%
  \newsavebox\ldv@publishersbox%
  \savebox\ldv@publishersbox[\textwidth]{%
    \parbox{\ldv@ldvlogowidth}{%
      \includegraphics[width=\ldv@ldvlogowidth]{\ldv@ldvlogoS}%
    }%
    \hfill%
    \parbox{\ldv@publisherwidth}{%
      \centering%
      \@institute\\%
      \@publishers%
    }%
    \hfill%
    \parbox{\ldv@tumlogowidth}{%
      \includegraphics[width=\ldv@tumlogowidth]{\ldv@tumlogoS}%
    }%
  }%
  %
  \newlength\ldv@bottommargin
  \setlength\ldv@bottommargin{1in + \hoffset + \oddsidemargin}%
  \newlength\ldv@publishersboxdepth%
  \settodepth\ldv@publishersboxdepth{\usebox\ldv@publishersbox}%
  \newlength\ldv@publishersbottom%
  \setlength\ldv@publishersbottom{%
    1in + \voffset + \topmargin + \headheight + \headsep +%
    \topskip + \ldv@publishersboxdepth%
  }%
  \newlength\ldv@publishersmove%
  \setlength\ldv@publishersmove{%
    \paperheight - \ldv@publishersbottom - \ldv@bottommargin%
  }%
  \noindent\raisebox{-\ldv@publishersmove}[0pt][0pt]{%
    \makebox[0pt][l]{\usebox\ldv@publishersbox}}%
}
\fi
%    \end{macrocode}
% \end{macro}
%
%
%
% \subsubsection{Titelblatt für Dissertationen}
% \label{sec:titelblatt-dissertationen}
%
% \begin{macro}{\ldv@phdtitle}
%    \begin{macrocode}
%<*book>
\newcommand*\@dateaccepted{}
\newcommand*\dateaccepted[1]{\renewcommand*\@dateaccepted{#1}}
\newcommand*\@pruefer{}
\newcommand*\pruefer[1]{\renewcommand*\@pruefer{#1}}
\newcommand*\@vorsitzender{}
\newcommand*\vorsitzender[1]{\renewcommand*\@vorsitzender{#1}}
\newcommand*\ldv@phdtitle{%
  \setparsizes{\z@}{\z@}{\z@\@plus 1fil}\par@updaterelative%
  \newsavebox\ldv@publishersbox%
  \savebox\ldv@publishersbox[\textwidth]{%
    \hfill%
    \parbox{\ldv@tumlogowidth}{%
      \includegraphics[width=\ldv@tumlogowidth]{\ldv@tumlogoS}%
    }%
  }%
  \newlength\ldv@topmargin%
  \setlength\ldv@topmargin{0mm}%
  \newlength\ldv@publishersboxheight%
  \settodepth\ldv@publishersboxheight{\usebox\ldv@publishersbox}%
  \newlength\ldv@publisherstop%
  \setlength\ldv@publisherstop{%
    1in + \voffset + \topmargin + \headheight + \headsep +%
    \topskip - \ldv@publishersboxheight%
  }%
  \newlength\ldv@publishersmove%
  % TODO Rewrite page instead of using magic number for
  % text/logo alignment
  \setlength\ldv@publishersmove{-3.5pt}%
  \noindent\raisebox{\ldv@publishersmove}[0mm][0mm]{%
    \makebox[0mm][l]{\usebox\ldv@publishersbox}}%
  Lehrstuhl f\"ur Datenverarbeitung\\
  Technische Universit\"at M\"unchen
  \vspace{3\baselineskip}

  % TODO Make this font adjustable with the KOMA-Script font
  % selection mechanism
  {\raggedright\LARGE\bfseries \ifx\@titletitle\@empty \@title
    \else\@titletitle \fi\par}%
  \ifx\@subtitle\@empty\else{\vspace{0.5\baselineskip}\raggedright\large\bfseries\@subtitle\par}\fi%
  \vspace{2\baselineskip}

  {\large\bfseries\@author\par}%
  \vspace{2\baselineskip}

  \begin{otherlanguage}{ngerman}% This text is always in German
    \hyphenation{In-for-ma-tions-tech-nik}%
    %
    \newlength\ldv@widthoftext%
    \settowidth\ldv@widthoftext{Vollst\"andiger Abdruck der von der
      Fakult\"at f\"ur Elektrotechnik und Informationstechnik}%
    \ifdim\ldv@widthoftext>\columnwidth%
      Vollst\"andiger Abdruck der von der Fakult\"at f\"ur
      Elektrotechnik und Informationstechnik der Technischen
      Universit\"at M\"unchen zur Erlangung des akademischen Grades
      eines\\[\baselineskip]%
    \else%
      Vollst\"andiger Abdruck der von der Fakult\"at f\"ur
      Elektrotechnik und Informationstechnik\\
      der Technischen Universit\"at M\"unchen zur Erlangung des
      akademischen Grades eines\\[\baselineskip]%
    \fi%
    {\bfseries Doktor-Ingenieurs (Dr.-Ing.)}\\[\baselineskip]%
    genehmigten Dissertation.  \vspace{3\baselineskip}%

    {\bfseries Vorsitzende(r):}\quad \@vorsitzender
    \vspace{\baselineskip}%

    {\bfseries Pr\"ufer der Dissertation:}%
    \let\tempand=\and%
    \renewcommand*\and{\item}%
    \begin{enumerate}%
    \item \@pruefer%
    \end{enumerate}%
    \let\and=\tempand%
    \vspace{.5\baselineskip}

    Die Dissertation wurde am \@date\ bei der Technischen
    Universit\"at M\"unchen eingereicht und durch die Fakult\"at f\"ur
    Elektrotechnik und Informationstechnik am \@dateaccepted\
    angenommen.
  \end{otherlanguage}%
}
%</book>
%    \end{macrocode}
% \end{macro}
%
%
%
% \subsubsection{Kleiner Titel am Seitenkopf}
%
% \begin{macro}{\ldv@headtitle}
%   Damit wird im Fall des Titelkopfes die eigentliche Arbeit
%   geleistet. Dieses Makro wurde aus den \KOMAScript-Quellen
%   übernommen, wo es \verb|\@maketitle| hieß, und an das TUM-Layout
%   angepasst. Deshalb kann das (ursprüngliche) Makro
%   \verb|@maketitle| gelöscht werden.
%
%   Falls auf ein anderes Papierformat gegangen wird, müsste
%   eigentlich auch die Logogröße angepasst werden (A5: 8mm, A4: 10mm,
%   A3: 14mm). Bis jetzt ist das noch nicht drinnen. Sollte noch
%   folgen. (!!)
%
%   Änderungen zu den \KOMAScript-Quellen:
%   \begin{itemize}
%   \item Logos eingefügt.
%   \item Verlag/Institution entfernt, weil er bereits durch die Logos
%     angezeigt ist.
%   \item Autor und Datum in normaler Schriftgröße (ohne
%     \verb|\Large|).
%   \item Versionsinformation beim Datum eingefügt.
%   \end{itemize}
%
%    \begin{macrocode}
\if@titlepage\else
\newcommand*\ldv@headtitle{%
  \clearpage
  \let\footnote\thanks%
  \ifx\@extratitle\@empty \else
    \noindent\@extratitle \next@tpage \if@twoside \null\next@tpage \fi
  \fi
  \noindent\includegraphics[width=\ldv@ldvlogowidth]{\ldv@ldvlogoS}%
  \hfill%
  \includegraphics[width=\ldv@tumlogowidth]{\ldv@tumlogoS}\par%
  \setparsizes{\z@}{\z@}{\z@\@plus 1fil}\par@updaterelative
  \vskip .5em
  \ifx\@titlehead\@empty \else
    \begin{minipage}[t]{\textwidth}
      \@titlehead
    \end{minipage}\par
  \fi
  \null
  \vspace\baselineskip%
  \begin{center}%
    \ifx\@subject\@empty \else%
      \vspace{-\baselineskip}{\subject@font\@subject}\par%
      \vspace\baselineskip%
    \fi%
    {\titlefont\huge\ifx\@titletitle\@empty\@title\else\@titletitle\fi\par}%
    \ifx\@subtitle\@empty\else%
      \vspace{.33\baselineskip}%
      {\usekomafont{subtitle}\@subtitle}\par%
    \fi%
    \vspace\baselineskip%
    \ifx\@author\@empty\vspace{.5em}\else%
      {\lineskip .5em%
        \begin{tabular}[t]{c}
          \@author
        \end{tabular}\par%
      }%
      \vspace{.5\baselineskip}%
    \fi%
    {\@date%
      \ifx\@version\@empty\else, \@version\fi \par}%
    {\ifx\ldv@licensetext\@empty\else\if@twoside\else%
      \vspace{.5\baselineskip}%
      \footnotesize\ldv@licensetext\par%
    \fi\fi}%
    \ifx\@dedication\@empty \else
      \vskip 2em
      {\Large \@dedication \par}
    \fi
  \end{center}%
  \par
  \vskip 2em
}
\fi
%    \end{macrocode}
% \end{macro}
%
%
%
% \subsubsection{Impressumsseite}
%
% Beim zweiseitigen Druck erscheinen die Details zum Buch oder zur
% Abschlussarbeit auf der Rückseite der Titelseite -- so wie es in
% Büchern üblich ist.
%
%    \begin{macrocode}
%<*book>
\if@twoside%
%    \end{macrocode}
%
% \begin{macro}{\makereference}
%   Zentraler Bestandteil des Impressums soll eine vollständige
%   Referenzierung des Werkes sein. Sie wird vom Makro
%   \verb|\makereference| zusammengesetzt.
%
%    \begin{macrocode}
  \newcommand*\makereference{
    \let\ldv@sep=\@empty
    \ifx\@author\@empty\else\authorwithand{, }. \fi%
    \ifx\@title\@empty\else\emph{\@title. \ifx\@subtitle\@empty\else\@subtitle. \fi}\fi%
    \ifx\@version\@empty\else\@version. \fi%
    \ifx\@subject\@empty%
    \else%
      \@subject%
      \renewcommand*\ldv@sep{, }%
    \fi%
    \ifx\@publishers\@empty%
    \else%
      \ldv@sep\@publishers, \@citationaddress%
      \renewcommand*\ldv@sep{, }%
    \fi
    \ldv@sep\number\year.%
  }
%    \end{macrocode}
% \end{macro}
%
% \begin{macro}{\ldv@thesissubmissiontext}
%     Das Impressum wird mit dem Makro \verb|\lowertitleback| der
%     \KOMAScript-Klassen gesetzt.  Der Inhalt hängt davon ab, ob es
%     sich um eine studentische Abschlussarbeit oder ein normales Buch
%     handelt.
%
%     Für den Fall einer studentischen Abschlussarbeit enthält das
%     Makro \verb|\ldv@thesissubmissiontext| den zusätzlichen
%     sprachenabhängigen Text.
%
%    \begin{macrocode}
  %
  \ifldv@studthesis%
    \newcommand*\ldv@thesissubmissiontext{%
      Supervised by Prof.\ Dr.-Ing.\ Klaus Diepold
      \ifx\@supervisor\@empty\else and \@supervisor\fi; submitted on
      \@date\ to the Department of Electrical and Computer
      Engineering of the \@publishers.%
    }
    %
    \addto\captionsenglish{%
      \renewcommand*\ldv@thesissubmissiontext{%
        Supervised by Prof.\ Dr.-Ing.\ Klaus Diepold
        \ifx\@supervisor\@empty\else and \@supervisor\fi; submitted on
        \@date\ to the Department of Electrical and Computer
        Engineering of the \@publishers.%
      }%
    }
    \addto\captionsngerman{%
      \renewcommand*\ldv@thesissubmissiontext{%
        Betreut von Prof.\ Dr.-Ing.\ Klaus Diepold
        \ifx\@supervisor\@empty\else und \@supervisor\fi; eingereicht
        am \@date\ bei der Fakult\"at f\"ur Elektrotechnik und
        Informationstechnik der Technischen Universit\"at M\"unchen.
      }%
    }
    \addto\captionsgerman{%
      \renewcommand*\ldv@thesissubmissiontext{%
        Betreut von Prof.\ Dr.-Ing.\ Klaus Diepold
        \ifx\@supervisor\@empty\else und \@supervisor\fi; eingereicht
        am \@date\ bei der Fakult\"at f\"ur Elektrotechnik und
        Informationstechnik der Technischen Universit\"at M\"unchen.
      }%
    }
    %
    \lowertitleback{%
      \makereference
      %
      \ifx\@keywords\@empty\else\\[\baselineskip]\keywordsname: \@keywords .\fi%
      %
      \\[\baselineskip]\ldv@thesissubmissiontext
      %
      \ifx\ldv@licensetext\@empty\else%
        \\[\baselineskip]
        \copyright\ \number\year\ \authorwithand{, }\\[\baselineskip]
        %
        \@institute, \@publishers, \@postaddress,
        \@publishersurl.\\[\baselineskip]
        %
        \ldv@licensetext%
      \fi%
    }%
  \else
    \lowertitleback{%
      \makereference%
      %
      \ifx\@keywords\@empty\else\\[\baselineskip]\keywordsname: \@keywords .\fi%
      %
      \ifx\ldv@licensetext\@empty\else
        \\[\baselineskip]
        \copyright\ \number\year\ \authorwithand{, }\\[\baselineskip]
        %
        \ifx\@publishers\@empty\else%
          \ifx\@institute\@empty\else\@institute, \fi%
          \@publishers%
          \ifx\@postaddress\@empty\else, \@postaddress\fi%
          \ifx\@publishersurl\@empty\else, \@publishersurl\fi%
          .\\[\baselineskip]%
        \fi%
        %
        \ldv@licensetext%
      \fi
    }%
  \fi
\fi
%</book>
%    \end{macrocode}
% \end{macro}
%
%
%
% \subsubsection{Umschlagseiten}
%
% Umschlagseiten bilden ein gestalterisch freies und individuelles
% Element an einem Buch. Um der Freiheit Rechnung zu tragen, soll es
% in einer späteren Version der Klassen die Möglichkeit geben, eine
% selbst gestaltete Seite im PDF-Format einzubinden. Darüber hinaus
% sollen die Klassen mehrere verschiedene vordefinierte Designs
% anbieten, aus denen sich der Autor dann eines aussuchen kann. Zur
% Zeit bieten die LDV-Klassen aber nur ein Umschlaglayout.
%
% \begin{macro}{\@covertitle}
% \begin{macro}{\covertitle}
%   Um für die Umschlagseite eigene Zeilenumbrüche im Titel zu setzen,
%   kann der Benutzer das Hilfsmakro \verb|\covertitle| benutzen. Der
%   so festgelegte Titel wird dann beim Setzen des Umschlags
%   bevorzugt.
%
%    \begin{macrocode}
\newcommand*\@covertitle{}
\newcommand*\covertitle[1]{%
  \renewcommand*\@covertitle{#1}%
}
%    \end{macrocode}
% \end{macro}
% \end{macro}
%
% \begin{macro}{\ldv@makecover@eins}
%   Das Makro \verb|\ldv@makecover@eins| setzt das \verb|Design1| der
%   vorderen Umschlagseite. Es setzt das Verhältnis des goldenen
%   Schnitts um: Gesamtbreite zu linker Spalte zu rechter Spalte.
%
%    \begin{macrocode}
\if@titlepage
\newcommand\ldv@makecover@eins{%
% Die verfügbare Spaltenbreite entspricht der eigentlichen Textbreite,
% d.h. die Ränder müssen abgezogen werden. Damit passt sich das Layout
% an unterschiedliche DIV und BCOR Werte automatisch an.
  \newlength\ldv@coverleftcolumnwidth%
  \setlength\ldv@coverleftcolumnwidth{0.618\paperwidth %
    -\hoffset %
    -1in %
    -\oddsidemargin %
    -2mm}%
% Für die rechte Spalte nutzen wir zusätzlich den Platz für Randnotizen,
% damit steht auch für sehr lange Namen genügend Platz zur Verfügung.
  \newlength\ldv@coverrightcolumnwidth%
  \setlength\ldv@coverrightcolumnwidth{0.382\paperwidth %
    -\hoffset %
    -1in %
    -\evensidemargin %
    -2mm %
    +\marginparsep %
    +\marginparwidth}%
  \newsavebox\ldv@covertitle%
  \savebox\ldv@covertitle{%
    \parbox{\ldv@coverleftcolumnwidth}{%
      \raggedleft\bfseries%
      {\huge\ifx\@covertitle\empty\@title\else\@covertitle\fi\medskip\par}
      \@subtitle%
    }%
  }%
  \newsavebox\ldv@coverauthor%
  \savebox\ldv@coverauthor{%
    \parbox{\ldv@coverrightcolumnwidth}{%
      \raggedright\bfseries%
      \authorwithand{\newline}%
    }%
  }%
% Der Abstand von der Oberkante entspricht dem Standardlayout
  \newlength\ldv@margintop%
  \setlength\ldv@margintop{1in %
    +\voffset %
    +\topmargin %
    +\headheight %
    +\headsep%
  }
  \newlength\ldv@authorpadding%
  \setlength\ldv@authorpadding{\ldv@margintop %
    +\topskip %
    +\totalheightof{\usebox\ldv@covertitle} %
    +1mm %
  }%
% Das eigentliche Layout wird mit TikZ gezeichnet, das ist zugänglicher und
% erzeugt keine "overfull hbox"-Warnungen.
  \begin{tikzpicture}[remember picture, overlay]
    \node [
      below left,
      align=right,
      xshift=\hoffset + 0.618\paperwidth - 2mm,
      yshift=-\ldv@margintop - \topskip,
    ] at (current page.north west) {\usebox\ldv@covertitle};
    \node [
      below right,
      align=left,
      xshift=\hoffset + 0.618\paperwidth + 2mm,
      yshift=-\ldv@authorpadding,
    ] at (current page.north west) {\usebox\ldv@coverauthor};
    \path[fill=TUMBlau5] (current page.north west) ++(\hoffset + 0.618\paperwidth - 1.5mm, 0) rectangle ++(3mm, -\paperheight);
    \node [
      below left,
      xshift=\hoffset + 0.618\paperwidth - 2mm,
      yshift=-\ldv@margintop - \textheight - \footskip,
    ] at (current page.north west) {\includegraphics[width=\ldv@ldvlogowidth]{\ldv@ldvlogoCMYK}};
    \node [
      below right,
      xshift=\hoffset + 0.618\paperwidth + 2mm,
      yshift=-\ldv@margintop - \textheight - \footskip,
    ] at (current page.north west) {\includegraphics[width=\ldv@tumlogowidth]{\ldv@tumlogoCMYK}};
  \end{tikzpicture}%
}%
\fi
%    \end{macrocode}
% \end{macro}
%
%
% \begin{macro}{\ldv@includefile}
%    \begin{macrocode}
\newcommand*\ldv@includefile[1]{%
  \noindent\hspace{-1in}\hspace{-\oddsidemargin}\fbox{A}%
  \raisebox{1in + \topmargin + \headheight + \headsep}[0mm][0mm]{%
    \fbox{B\includegraphics{LDVLogoS_oT}}
  }
%  \noindent\hspace{-1in}\hspace{-\oddsidemargin}%
%  \raisebox{1in + \topmargin + \headheight + \headsep +
%    \topskip}{%
%    \fbox{\includegraphics{ta.pdf}}%
%  }
}
%    \end{macrocode}
% \end{macro}
%
%
%
\endinput
%
%
%
% \iffalse
%%% Local Variables:
%%% mode: doctex
%%% TeX-PDF-mode: t
%%% TeX-master: "ldvguide.tex"
%%% End:
% \fi
